% =====================================================================
%  Taller – Análisis de Componentes Principales (ACP) – Base Wine (UCI)
%  Compilar en Overleaf con pdfLaTeX. Todos los gráficos están hechos
%  con pgfplots y los datos van incluidos en este mismo archivo.
% =====================================================================
\documentclass[11pt,a4paper]{article}

\usepackage[utf8]{inputenc}
\usepackage[T1]{fontenc}
\usepackage[spanish,es-nodecimaldot,es-noshorthands,es-tabla]{babel}
\usepackage{lmodern}
\usepackage[margin=2.3cm]{geometry}
\usepackage{amsmath,amssymb}
\usepackage{booktabs,multirow,array,tabularx}
\usepackage{xcolor}
\usepackage{graphicx}
\usepackage{float}
\usepackage{caption}
\usepackage{subcaption}
\usepackage{enumitem}
\usepackage{csquotes}
\usepackage[hidelinks]{hyperref}

\usepackage{pgfplots}
\usepackage{pgfplotstable}
\pgfplotsset{compat=1.18}
\usepgfplotslibrary{statistics,groupplots,colormaps}
\usetikzlibrary{arrows.meta,babel}

% ---------- colores ----------
\definecolor{c1}{HTML}{1B9E77}   % Cultivar 1
\definecolor{c2}{HTML}{D95F02}   % Cultivar 2
\definecolor{c3}{HTML}{7570B3}   % Cultivar 3
\definecolor{azul}{HTML}{4682B4}
\definecolor{caja}{HTML}{B0C4DE}
\definecolor{rojo}{HTML}{C0392B}
\definecolor{azulcos}{HTML}{1F3A93}

% filtro para graficar por cultivar
\pgfplotsset{
  discard if not/.style 2 args={restrict expr to domain={\thisrow{#1}}{#2:#2}},
  ejes/.style={axis lines=box, grid=major, grid style={gray!20},
    tick label style={font=\footnotesize}, label style={font=\small},
    title style={font=\small\bfseries}},
  cero/.style={gray, dashed, forget plot}
}

\newcommand{\var}[1]{\texttt{#1}}
\setlength{\parskip}{4pt}

% =====================================================================
%  DATOS (178 vinos)
% =====================================================================
\pgfplotstableread{
id alcohol acmalico ceniza alcalceniza magnesio fenolestot flavonoides fenolesnoflav proanto intcolor tono odratio prolina cult
1 14.23 1.71 2.43 15.6 127 2.8 3.06 0.28 2.29 5.64 1.04 3.92 1065 1
2 13.2 1.78 2.14 11.2 100 2.65 2.76 0.26 1.28 4.38 1.05 3.4 1050 1
3 13.16 2.36 2.67 18.6 101 2.8 3.24 0.3 2.81 5.68 1.03 3.17 1185 1
4 14.37 1.95 2.5 16.8 113 3.85 3.49 0.24 2.18 7.8 0.86 3.45 1480 1
5 13.24 2.59 2.87 21 118 2.8 2.69 0.39 1.82 4.32 1.04 2.93 735 1
6 14.2 1.76 2.45 15.2 112 3.27 3.39 0.34 1.97 6.75 1.05 2.85 1450 1
7 14.39 1.87 2.45 14.6 96 2.5 2.52 0.3 1.98 5.25 1.02 3.58 1290 1
8 14.06 2.15 2.61 17.6 121 2.6 2.51 0.31 1.25 5.05 1.06 3.58 1295 1
9 14.83 1.64 2.17 14 97 2.8 2.98 0.29 1.98 5.2 1.08 2.85 1045 1
10 13.86 1.35 2.27 16 98 2.98 3.15 0.22 1.85 7.22 1.01 3.55 1045 1
11 14.1 2.16 2.3 18 105 2.95 3.32 0.22 2.38 5.75 1.25 3.17 1510 1
12 14.12 1.48 2.32 16.8 95 2.2 2.43 0.26 1.57 5 1.17 2.82 1280 1
13 13.75 1.73 2.41 16 89 2.6 2.76 0.29 1.81 5.6 1.15 2.9 1320 1
14 14.75 1.73 2.39 11.4 91 3.1 3.69 0.43 2.81 5.4 1.25 2.73 1150 1
15 14.38 1.87 2.38 12 102 3.3 3.64 0.29 2.96 7.5 1.2 3 1547 1
16 13.63 1.81 2.7 17.2 112 2.85 2.91 0.3 1.46 7.3 1.28 2.88 1310 1
17 14.3 1.92 2.72 20 120 2.8 3.14 0.33 1.97 6.2 1.07 2.65 1280 1
18 13.83 1.57 2.62 20 115 2.95 3.4 0.4 1.72 6.6 1.13 2.57 1130 1
19 14.19 1.59 2.48 16.5 108 3.3 3.93 0.32 1.86 8.7 1.23 2.82 1680 1
20 13.64 3.1 2.56 15.2 116 2.7 3.03 0.17 1.66 5.1 0.96 3.36 845 1
21 14.06 1.63 2.28 16 126 3 3.17 0.24 2.1 5.65 1.09 3.71 780 1
22 12.93 3.8 2.65 18.6 102 2.41 2.41 0.25 1.98 4.5 1.03 3.52 770 1
23 13.71 1.86 2.36 16.6 101 2.61 2.88 0.27 1.69 3.8 1.11 4 1035 1
24 12.85 1.6 2.52 17.8 95 2.48 2.37 0.26 1.46 3.93 1.09 3.63 1015 1
25 13.5 1.81 2.61 20 96 2.53 2.61 0.28 1.66 3.52 1.12 3.82 845 1
26 13.05 2.05 3.22 25 124 2.63 2.68 0.47 1.92 3.58 1.13 3.2 830 1
27 13.39 1.77 2.62 16.1 93 2.85 2.94 0.34 1.45 4.8 0.92 3.22 1195 1
28 13.3 1.72 2.14 17 94 2.4 2.19 0.27 1.35 3.95 1.02 2.77 1285 1
29 13.87 1.9 2.8 19.4 107 2.95 2.97 0.37 1.76 4.5 1.25 3.4 915 1
30 14.02 1.68 2.21 16 96 2.65 2.33 0.26 1.98 4.7 1.04 3.59 1035 1
31 13.73 1.5 2.7 22.5 101 3 3.25 0.29 2.38 5.7 1.19 2.71 1285 1
32 13.58 1.66 2.36 19.1 106 2.86 3.19 0.22 1.95 6.9 1.09 2.88 1515 1
33 13.68 1.83 2.36 17.2 104 2.42 2.69 0.42 1.97 3.84 1.23 2.87 990 1
34 13.76 1.53 2.7 19.5 132 2.95 2.74 0.5 1.35 5.4 1.25 3 1235 1
35 13.51 1.8 2.65 19 110 2.35 2.53 0.29 1.54 4.2 1.1 2.87 1095 1
36 13.48 1.81 2.41 20.5 100 2.7 2.98 0.26 1.86 5.1 1.04 3.47 920 1
37 13.28 1.64 2.84 15.5 110 2.6 2.68 0.34 1.36 4.6 1.09 2.78 880 1
38 13.05 1.65 2.55 18 98 2.45 2.43 0.29 1.44 4.25 1.12 2.51 1105 1
39 13.07 1.5 2.1 15.5 98 2.4 2.64 0.28 1.37 3.7 1.18 2.69 1020 1
40 14.22 3.99 2.51 13.2 128 3 3.04 0.2 2.08 5.1 0.89 3.53 760 1
41 13.56 1.71 2.31 16.2 117 3.15 3.29 0.34 2.34 6.13 0.95 3.38 795 1
42 13.41 3.84 2.12 18.8 90 2.45 2.68 0.27 1.48 4.28 0.91 3 1035 1
43 13.88 1.89 2.59 15 101 3.25 3.56 0.17 1.7 5.43 0.88 3.56 1095 1
44 13.24 3.98 2.29 17.5 103 2.64 2.63 0.32 1.66 4.36 0.82 3 680 1
45 13.05 1.77 2.1 17 107 3 3 0.28 2.03 5.04 0.88 3.35 885 1
46 14.21 4.04 2.44 18.9 111 2.85 2.65 0.3 1.25 5.24 0.87 3.33 1080 1
47 14.38 3.59 2.28 16 102 3.25 3.17 0.27 2.19 4.9 1.04 3.44 1065 1
48 13.9 1.68 2.12 16 101 3.1 3.39 0.21 2.14 6.1 0.91 3.33 985 1
49 14.1 2.02 2.4 18.8 103 2.75 2.92 0.32 2.38 6.2 1.07 2.75 1060 1
50 13.94 1.73 2.27 17.4 108 2.88 3.54 0.32 2.08 8.9 1.12 3.1 1260 1
51 13.05 1.73 2.04 12.4 92 2.72 3.27 0.17 2.91 7.2 1.12 2.91 1150 1
52 13.83 1.65 2.6 17.2 94 2.45 2.99 0.22 2.29 5.6 1.24 3.37 1265 1
53 13.82 1.75 2.42 14 111 3.88 3.74 0.32 1.87 7.05 1.01 3.26 1190 1
54 13.77 1.9 2.68 17.1 115 3 2.79 0.39 1.68 6.3 1.13 2.93 1375 1
55 13.74 1.67 2.25 16.4 118 2.6 2.9 0.21 1.62 5.85 0.92 3.2 1060 1
56 13.56 1.73 2.46 20.5 116 2.96 2.78 0.2 2.45 6.25 0.98 3.03 1120 1
57 14.22 1.7 2.3 16.3 118 3.2 3 0.26 2.03 6.38 0.94 3.31 970 1
58 13.29 1.97 2.68 16.8 102 3 3.23 0.31 1.66 6 1.07 2.84 1270 1
59 13.72 1.43 2.5 16.7 108 3.4 3.67 0.19 2.04 6.8 0.89 2.87 1285 1
60 12.37 0.94 1.36 10.6 88 1.98 0.57 0.28 0.42 1.95 1.05 1.82 520 2
61 12.33 1.1 2.28 16 101 2.05 1.09 0.63 0.41 3.27 1.25 1.67 680 2
62 12.64 1.36 2.02 16.8 100 2.02 1.41 0.53 0.62 5.75 0.98 1.59 450 2
63 13.67 1.25 1.92 18 94 2.1 1.79 0.32 0.73 3.8 1.23 2.46 630 2
64 12.37 1.13 2.16 19 87 3.5 3.1 0.19 1.87 4.45 1.22 2.87 420 2
65 12.17 1.45 2.53 19 104 1.89 1.75 0.45 1.03 2.95 1.45 2.23 355 2
66 12.37 1.21 2.56 18.1 98 2.42 2.65 0.37 2.08 4.6 1.19 2.3 678 2
67 13.11 1.01 1.7 15 78 2.98 3.18 0.26 2.28 5.3 1.12 3.18 502 2
68 12.37 1.17 1.92 19.6 78 2.11 2 0.27 1.04 4.68 1.12 3.48 510 2
69 13.34 0.94 2.36 17 110 2.53 1.3 0.55 0.42 3.17 1.02 1.93 750 2
70 12.21 1.19 1.75 16.8 151 1.85 1.28 0.14 2.5 2.85 1.28 3.07 718 2
71 12.29 1.61 2.21 20.4 103 1.1 1.02 0.37 1.46 3.05 0.906 1.82 870 2
72 13.86 1.51 2.67 25 86 2.95 2.86 0.21 1.87 3.38 1.36 3.16 410 2
73 13.49 1.66 2.24 24 87 1.88 1.84 0.27 1.03 3.74 0.98 2.78 472 2
74 12.99 1.67 2.6 30 139 3.3 2.89 0.21 1.96 3.35 1.31 3.5 985 2
75 11.96 1.09 2.3 21 101 3.38 2.14 0.13 1.65 3.21 0.99 3.13 886 2
76 11.66 1.88 1.92 16 97 1.61 1.57 0.34 1.15 3.8 1.23 2.14 428 2
77 13.03 0.9 1.71 16 86 1.95 2.03 0.24 1.46 4.6 1.19 2.48 392 2
78 11.84 2.89 2.23 18 112 1.72 1.32 0.43 0.95 2.65 0.96 2.52 500 2
79 12.33 0.99 1.95 14.8 136 1.9 1.85 0.35 2.76 3.4 1.06 2.31 750 2
80 12.7 3.87 2.4 23 101 2.83 2.55 0.43 1.95 2.57 1.19 3.13 463 2
81 12 0.92 2 19 86 2.42 2.26 0.3 1.43 2.5 1.38 3.12 278 2
82 12.72 1.81 2.2 18.8 86 2.2 2.53 0.26 1.77 3.9 1.16 3.14 714 2
83 12.08 1.13 2.51 24 78 2 1.58 0.4 1.4 2.2 1.31 2.72 630 2
84 13.05 3.86 2.32 22.5 85 1.65 1.59 0.61 1.62 4.8 0.84 2.01 515 2
85 11.84 0.89 2.58 18 94 2.2 2.21 0.22 2.35 3.05 0.79 3.08 520 2
86 12.67 0.98 2.24 18 99 2.2 1.94 0.3 1.46 2.62 1.23 3.16 450 2
87 12.16 1.61 2.31 22.8 90 1.78 1.69 0.43 1.56 2.45 1.33 2.26 495 2
88 11.65 1.67 2.62 26 88 1.92 1.61 0.4 1.34 2.6 1.36 3.21 562 2
89 11.64 2.06 2.46 21.6 84 1.95 1.69 0.48 1.35 2.8 1 2.75 680 2
90 12.08 1.33 2.3 23.6 70 2.2 1.59 0.42 1.38 1.74 1.07 3.21 625 2
91 12.08 1.83 2.32 18.5 81 1.6 1.5 0.52 1.64 2.4 1.08 2.27 480 2
92 12 1.51 2.42 22 86 1.45 1.25 0.5 1.63 3.6 1.05 2.65 450 2
93 12.69 1.53 2.26 20.7 80 1.38 1.46 0.58 1.62 3.05 0.96 2.06 495 2
94 12.29 2.83 2.22 18 88 2.45 2.25 0.25 1.99 2.15 1.15 3.3 290 2
95 11.62 1.99 2.28 18 98 3.02 2.26 0.17 1.35 3.25 1.16 2.96 345 2
96 12.47 1.52 2.2 19 162 2.5 2.27 0.32 3.28 2.6 1.16 2.63 937 2
97 11.81 2.12 2.74 21.5 134 1.6 0.99 0.14 1.56 2.5 0.95 2.26 625 2
98 12.29 1.41 1.98 16 85 2.55 2.5 0.29 1.77 2.9 1.23 2.74 428 2
99 12.37 1.07 2.1 18.5 88 3.52 3.75 0.24 1.95 4.5 1.04 2.77 660 2
100 12.29 3.17 2.21 18 88 2.85 2.99 0.45 2.81 2.3 1.42 2.83 406 2
101 12.08 2.08 1.7 17.5 97 2.23 2.17 0.26 1.4 3.3 1.27 2.96 710 2
102 12.6 1.34 1.9 18.5 88 1.45 1.36 0.29 1.35 2.45 1.04 2.77 562 2
103 12.34 2.45 2.46 21 98 2.56 2.11 0.34 1.31 2.8 0.8 3.38 438 2
104 11.82 1.72 1.88 19.5 86 2.5 1.64 0.37 1.42 2.06 0.94 2.44 415 2
105 12.51 1.73 1.98 20.5 85 2.2 1.92 0.32 1.48 2.94 1.04 3.57 672 2
106 12.42 2.55 2.27 22 90 1.68 1.84 0.66 1.42 2.7 0.86 3.3 315 2
107 12.25 1.73 2.12 19 80 1.65 2.03 0.37 1.63 3.4 1 3.17 510 2
108 12.72 1.75 2.28 22.5 84 1.38 1.76 0.48 1.63 3.3 0.88 2.42 488 2
109 12.22 1.29 1.94 19 92 2.36 2.04 0.39 2.08 2.7 0.86 3.02 312 2
110 11.61 1.35 2.7 20 94 2.74 2.92 0.29 2.49 2.65 0.96 3.26 680 2
111 11.46 3.74 1.82 19.5 107 3.18 2.58 0.24 3.58 2.9 0.75 2.81 562 2
112 12.52 2.43 2.17 21 88 2.55 2.27 0.26 1.22 2 0.9 2.78 325 2
113 11.76 2.68 2.92 20 103 1.75 2.03 0.6 1.05 3.8 1.23 2.5 607 2
114 11.41 0.74 2.5 21 88 2.48 2.01 0.42 1.44 3.08 1.1 2.31 434 2
115 12.08 1.39 2.5 22.5 84 2.56 2.29 0.43 1.04 2.9 0.93 3.19 385 2
116 11.03 1.51 2.2 21.5 85 2.46 2.17 0.52 2.01 1.9 1.71 2.87 407 2
117 11.82 1.47 1.99 20.8 86 1.98 1.6 0.3 1.53 1.95 0.95 3.33 495 2
118 12.42 1.61 2.19 22.5 108 2 2.09 0.34 1.61 2.06 1.06 2.96 345 2
119 12.77 3.43 1.98 16 80 1.63 1.25 0.43 0.83 3.4 0.7 2.12 372 2
120 12 3.43 2 19 87 2 1.64 0.37 1.87 1.28 0.93 3.05 564 2
121 11.45 2.4 2.42 20 96 2.9 2.79 0.32 1.83 3.25 0.8 3.39 625 2
122 11.56 2.05 3.23 28.5 119 3.18 5.08 0.47 1.87 6 0.93 3.69 465 2
123 12.42 4.43 2.73 26.5 102 2.2 2.13 0.43 1.71 2.08 0.92 3.12 365 2
124 13.05 5.8 2.13 21.5 86 2.62 2.65 0.3 2.01 2.6 0.73 3.1 380 2
125 11.87 4.31 2.39 21 82 2.86 3.03 0.21 2.91 2.8 0.75 3.64 380 2
126 12.07 2.16 2.17 21 85 2.6 2.65 0.37 1.35 2.76 0.86 3.28 378 2
127 12.43 1.53 2.29 21.5 86 2.74 3.15 0.39 1.77 3.94 0.69 2.84 352 2
128 11.79 2.13 2.78 28.5 92 2.13 2.24 0.58 1.76 3 0.97 2.44 466 2
129 12.37 1.63 2.3 24.5 88 2.22 2.45 0.4 1.9 2.12 0.89 2.78 342 2
130 12.04 4.3 2.38 22 80 2.1 1.75 0.42 1.35 2.6 0.79 2.57 580 2
131 12.86 1.35 2.32 18 122 1.51 1.25 0.21 0.94 4.1 0.76 1.29 630 3
132 12.88 2.99 2.4 20 104 1.3 1.22 0.24 0.83 5.4 0.74 1.42 530 3
133 12.81 2.31 2.4 24 98 1.15 1.09 0.27 0.83 5.7 0.66 1.36 560 3
134 12.7 3.55 2.36 21.5 106 1.7 1.2 0.17 0.84 5 0.78 1.29 600 3
135 12.51 1.24 2.25 17.5 85 2 0.58 0.6 1.25 5.45 0.75 1.51 650 3
136 12.6 2.46 2.2 18.5 94 1.62 0.66 0.63 0.94 7.1 0.73 1.58 695 3
137 12.25 4.72 2.54 21 89 1.38 0.47 0.53 0.8 3.85 0.75 1.27 720 3
138 12.53 5.51 2.64 25 96 1.79 0.6 0.63 1.1 5 0.82 1.69 515 3
139 13.49 3.59 2.19 19.5 88 1.62 0.48 0.58 0.88 5.7 0.81 1.82 580 3
140 12.84 2.96 2.61 24 101 2.32 0.6 0.53 0.81 4.92 0.89 2.15 590 3
141 12.93 2.81 2.7 21 96 1.54 0.5 0.53 0.75 4.6 0.77 2.31 600 3
142 13.36 2.56 2.35 20 89 1.4 0.5 0.37 0.64 5.6 0.7 2.47 780 3
143 13.52 3.17 2.72 23.5 97 1.55 0.52 0.5 0.55 4.35 0.89 2.06 520 3
144 13.62 4.95 2.35 20 92 2 0.8 0.47 1.02 4.4 0.91 2.05 550 3
145 12.25 3.88 2.2 18.5 112 1.38 0.78 0.29 1.14 8.21 0.65 2 855 3
146 13.16 3.57 2.15 21 102 1.5 0.55 0.43 1.3 4 0.6 1.68 830 3
147 13.88 5.04 2.23 20 80 0.98 0.34 0.4 0.68 4.9 0.58 1.33 415 3
148 12.87 4.61 2.48 21.5 86 1.7 0.65 0.47 0.86 7.65 0.54 1.86 625 3
149 13.32 3.24 2.38 21.5 92 1.93 0.76 0.45 1.25 8.42 0.55 1.62 650 3
150 13.08 3.9 2.36 21.5 113 1.41 1.39 0.34 1.14 9.4 0.57 1.33 550 3
151 13.5 3.12 2.62 24 123 1.4 1.57 0.22 1.25 8.6 0.59 1.3 500 3
152 12.79 2.67 2.48 22 112 1.48 1.36 0.24 1.26 10.8 0.48 1.47 480 3
153 13.11 1.9 2.75 25.5 116 2.2 1.28 0.26 1.56 7.1 0.61 1.33 425 3
154 13.23 3.3 2.28 18.5 98 1.8 0.83 0.61 1.87 10.52 0.56 1.51 675 3
155 12.58 1.29 2.1 20 103 1.48 0.58 0.53 1.4 7.6 0.58 1.55 640 3
156 13.17 5.19 2.32 22 93 1.74 0.63 0.61 1.55 7.9 0.6 1.48 725 3
157 13.84 4.12 2.38 19.5 89 1.8 0.83 0.48 1.56 9.01 0.57 1.64 480 3
158 12.45 3.03 2.64 27 97 1.9 0.58 0.63 1.14 7.5 0.67 1.73 880 3
159 14.34 1.68 2.7 25 98 2.8 1.31 0.53 2.7 13 0.57 1.96 660 3
160 13.48 1.67 2.64 22.5 89 2.6 1.1 0.52 2.29 11.75 0.57 1.78 620 3
161 12.36 3.83 2.38 21 88 2.3 0.92 0.5 1.04 7.65 0.56 1.58 520 3
162 13.69 3.26 2.54 20 107 1.83 0.56 0.5 0.8 5.88 0.96 1.82 680 3
163 12.85 3.27 2.58 22 106 1.65 0.6 0.6 0.96 5.58 0.87 2.11 570 3
164 12.96 3.45 2.35 18.5 106 1.39 0.7 0.4 0.94 5.28 0.68 1.75 675 3
165 13.78 2.76 2.3 22 90 1.35 0.68 0.41 1.03 9.58 0.7 1.68 615 3
166 13.73 4.36 2.26 22.5 88 1.28 0.47 0.52 1.15 6.62 0.78 1.75 520 3
167 13.45 3.7 2.6 23 111 1.7 0.92 0.43 1.46 10.68 0.85 1.56 695 3
168 12.82 3.37 2.3 19.5 88 1.48 0.66 0.4 0.97 10.26 0.72 1.75 685 3
169 13.58 2.58 2.69 24.5 105 1.55 0.84 0.39 1.54 8.66 0.74 1.8 750 3
170 13.4 4.6 2.86 25 112 1.98 0.96 0.27 1.11 8.5 0.67 1.92 630 3
171 12.2 3.03 2.32 19 96 1.25 0.49 0.4 0.73 5.5 0.66 1.83 510 3
172 12.77 2.39 2.28 19.5 86 1.39 0.51 0.48 0.64 9.9 0.57 1.63 470 3
173 14.16 2.51 2.48 20 91 1.68 0.7 0.44 1.24 9.7 0.62 1.71 660 3
174 13.71 5.65 2.45 20.5 95 1.68 0.61 0.52 1.06 7.7 0.64 1.74 740 3
175 13.4 3.91 2.48 23 102 1.8 0.75 0.43 1.41 7.3 0.7 1.56 750 3
176 13.27 4.28 2.26 20 120 1.59 0.69 0.43 1.35 10.2 0.59 1.56 835 3
177 13.17 2.59 2.37 20 120 1.65 0.68 0.53 1.46 9.3 0.6 1.62 840 3
178 14.13 4.1 2.74 24.5 96 2.05 0.76 0.56 1.35 9.2 0.61 1.6 560 3
}\datos

\pgfplotstableread{
obs cpa cpb cpc cult
1 -0.2486 -0.1082 -0.0124 1
2 -0.1656 0.0250 -0.1519 1
3 -0.1886 -0.0773 0.0737 1
4 -0.2816 -0.2066 -0.0132 1
5 -0.0756 -0.0652 0.1519 1
6 -0.2286 -0.1591 -0.0472 1
7 -0.1836 -0.0881 -0.0732 1
8 -0.1544 -0.1206 0.0110 1
9 -0.1882 -0.0688 -0.1327 1
10 -0.2064 -0.0592 -0.0738 1
11 -0.2608 -0.0976 -0.0317 1
12 -0.1315 -0.0459 -0.0893 1
13 -0.1584 -0.0506 -0.0648 1
14 -0.2592 -0.0847 -0.0903 1
15 -0.3233 -0.1571 -0.0947 1
16 -0.1728 -0.1246 0.0163 1
17 -0.1628 -0.1744 0.0623 1
18 -0.1423 -0.1223 0.0596 1
19 -0.2655 -0.1888 -0.0364 1
20 -0.1562 -0.0795 -0.0123 1
21 -0.2342 -0.0590 -0.0273 1
22 -0.0814 -0.0181 0.0702 1
23 -0.1900 0.0069 -0.0234 1
24 -0.1233 0.0387 0.0108 1
25 -0.1320 0.0238 0.0667 1
26 -0.0742 -0.0705 0.2864 1
27 -0.1331 -0.0514 -0.0065 1
28 -0.0926 0.0067 -0.1040 1
29 -0.1640 -0.0517 0.1045 1
30 -0.1691 -0.0144 -0.0819 1
31 -0.1874 -0.0930 0.1039 1
32 -0.2007 -0.1103 -0.0249 1
33 -0.1221 -0.0040 -0.0125 1
34 -0.1426 -0.1224 0.0879 1
35 -0.1057 -0.0523 0.0360 1
36 -0.1427 -0.0132 0.0338 1
37 -0.1038 -0.0494 0.0344 1
38 -0.0841 -0.0086 -0.0029 1
39 -0.1126 0.0577 -0.1069 1
40 -0.1896 -0.1351 -0.0257 1
41 -0.1940 -0.0584 -0.0089 1
42 -0.0501 -0.0127 -0.0587 1
43 -0.2302 -0.0866 -0.0234 1
44 -0.0346 -0.0248 -0.0151 1
45 -0.1575 0.0053 -0.0492 1
46 -0.0852 -0.1332 0.0022 1
47 -0.2044 -0.0893 -0.0405 1
48 -0.2115 -0.0484 -0.0866 1
49 -0.1506 -0.0935 -0.0043 1
50 -0.2029 -0.1313 -0.0482 1
51 -0.2410 -0.0125 -0.1479 1
52 -0.2143 -0.0559 0.0004 1
53 -0.2628 -0.1209 -0.0390 1
54 -0.1668 -0.1405 0.0255 1
55 -0.1609 -0.0762 -0.0718 1
56 -0.1851 -0.0996 0.0385 1
57 -0.2055 -0.1077 -0.0459 1
58 -0.1629 -0.0909 0.0196 1
59 -0.2353 -0.1298 -0.0214 1
60 0.0696 0.2304 -0.3437 2
61 0.1156 0.1035 -0.0656 2
62 0.1376 0.0622 -0.1204 2
63 -0.0023 0.0946 -0.1337 2
64 -0.1537 0.1443 -0.0006 2
65 0.0457 0.1430 0.0509 2
66 -0.0675 0.0573 0.0430 2
67 -0.1685 0.1413 -0.1523 2
68 -0.0137 0.1819 -0.0802 2
69 0.0609 0.0165 -0.0530 2
70 -0.1481 0.1052 -0.0928 2
71 0.1178 0.0663 -0.0471 2
72 -0.1242 0.0717 0.1464 2
73 0.0544 0.0797 0.0060 2
74 -0.1920 -0.0195 0.2529 2
75 -0.1374 0.0965 0.0343 2
76 0.0651 0.1832 -0.1172 2
77 -0.0277 0.1614 -0.1836 2
78 0.1092 0.1037 -0.0170 2
79 -0.0947 0.0578 -0.0888 2
80 -0.0282 0.0770 0.1345 2
81 -0.0571 0.2530 -0.0268 2
82 -0.0775 0.1087 -0.0272 2
83 0.0371 0.1785 0.1001 2
84 0.1903 0.0066 0.0355 2
85 -0.0626 0.1105 0.0457 2
86 -0.0591 0.1519 -0.0191 2
87 0.0605 0.1674 0.0579 2
88 0.0418 0.1779 0.1730 2
89 0.0836 0.1351 0.0719 2
90 0.0417 0.1992 0.0636 2
91 0.1011 0.1588 -0.0036 2
92 0.1173 0.1388 0.0585 2
93 0.1449 0.1169 -0.0067 2
94 -0.0560 0.1734 0.0086 2
95 -0.0718 0.1667 0.0107 2
96 -0.1907 -0.0127 0.0591 2
97 0.0408 0.0277 0.0981 2
98 -0.0773 0.1923 -0.0814 2
99 -0.1688 0.1074 -0.0173 2
100 -0.1057 0.1624 0.0561 2
101 -0.0598 0.1782 -0.1175 2
102 0.0412 0.1719 -0.1123 2
103 0.0121 0.0873 0.0752 2
104 0.0495 0.2009 -0.0573 2
105 -0.0294 0.1573 -0.0354 2
106 0.1329 0.1287 0.0710 2
107 0.0275 0.1626 -0.0361 2
108 0.1215 0.1016 0.0215 2
109 -0.0062 0.1729 -0.0347 2
110 -0.1183 0.1096 0.1334 2
111 -0.1065 0.1063 0.0104 2
112 0.0209 0.1447 0.0059 2
113 0.0977 0.0572 0.1499 2
114 0.0343 0.1701 0.0796 2
115 0.0370 0.1453 0.0992 2
116 -0.0361 0.2902 0.1008 2
117 0.0190 0.2115 -0.0227 2
118 0.0080 0.1446 0.0517 2
119 0.1824 0.0942 -0.1426 2
120 0.0413 0.1666 -0.0267 2
121 -0.0554 0.1056 0.0843 2
122 -0.1002 -0.0190 0.4007 2
123 0.0882 0.0498 0.2256 2
124 0.0347 0.0463 0.0362 2
125 -0.0733 0.1084 0.1110 2
126 0.0073 0.1582 0.0326 2
127 -0.0029 0.0949 0.0515 2
128 0.1197 0.0906 0.2519 2
129 0.0359 0.1453 0.0972 2
130 0.1344 0.0862 0.0587 2
131 0.0995 -0.0128 -0.0884 3
132 0.1787 -0.0281 -0.0543 3
133 0.2201 -0.0198 -0.0126 3
134 0.1609 -0.0276 -0.0340 3
135 0.1776 0.0345 -0.0826 3
136 0.2296 -0.0265 -0.0824 3
137 0.2935 -0.0116 0.0166 3
138 0.2951 -0.0494 0.1283 3
139 0.2319 -0.0261 -0.0770 3
140 0.1780 -0.0219 0.0931 3
141 0.2083 -0.0215 0.0457 3
142 0.1714 -0.0279 -0.0728 3
143 0.2238 -0.0367 0.0710 3
144 0.1780 -0.0362 -0.0190 3
145 0.1656 -0.0869 -0.0933 3
146 0.1968 -0.0422 -0.0642 3
147 0.3208 -0.0487 -0.1093 3
148 0.2685 -0.0954 -0.0083 3
149 0.2104 -0.1177 -0.0354 3
150 0.2173 -0.1530 -0.0372 3
151 0.1739 -0.1766 0.0328 3
152 0.1911 -0.1533 -0.0234 3
153 0.1359 -0.1145 0.1021 3
154 0.2069 -0.1603 -0.0723 3
155 0.2052 -0.0307 -0.0892 3
156 0.2702 -0.1351 -0.0070 3
157 0.2166 -0.1443 -0.0586 3
158 0.2543 -0.0983 0.1201 3
159 0.0786 -0.2635 0.0869 3
160 0.1207 -0.1804 0.0411 3
161 0.2356 -0.0553 -0.0068 3
162 0.1679 -0.0881 -0.0076 3
163 0.2134 -0.0417 0.0603 3
164 0.1947 -0.0523 -0.0663 3
165 0.2211 -0.1166 -0.0737 3
166 0.2646 -0.0661 -0.0349 3
167 0.1803 -0.1943 0.0321 3
168 0.2195 -0.0955 -0.0909 3
169 0.1635 -0.1557 0.0572 3
170 0.1785 -0.1940 0.1063 3
171 0.2407 0.0188 -0.0635 3
172 0.2757 -0.0635 -0.1004 3
173 0.1848 -0.1644 -0.0689 3
174 0.2526 -0.1661 -0.0257 3
175 0.1950 -0.1317 0.0156 3
176 0.2007 -0.2069 -0.0705 3
177 0.1789 -0.1722 -0.0413 3
178 0.2405 -0.2075 0.0760 3
}\tablaF

\pgfplotstableread{
obs d1 d2
1 3.1732 1.7655
2 3.8996 0.8882
3 2.4261 0.0045
4 6.0297 1.6583
5 3.3409 0.3372
6 4.0323 1.0509
7 2.4736 0.8280
8 1.9419 0.9660
9 3.8245 2.5718
10 2.5166 1.8018
11 3.3570 0.1835
12 1.7750 0.3518
13 1.6403 0.0477
14 4.0333 4.0050
15 6.7783 0.2065
16 2.2563 1.5019
17 3.6295 2.5773
18 2.2564 3.4190
19 5.3387 2.7072
20 1.3853 0.5214
21 2.4009 1.3282
22 0.8764 3.6387
23 1.4284 1.3888
24 0.6922 2.5160
25 1.2408 1.0792
26 10.5987 0.1031
27 0.8586 0.2570
28 1.6484 0.8391
29 2.5387 0.3942
30 1.9111 2.0271
31 3.2550 2.2589
32 2.4534 0.7146
33 0.5808 0.8265
34 2.7717 0.0983
35 0.7726 0.3943
36 0.9181 0.4892
37 0.7225 0.1638
38 0.2723 0.9536
39 2.1113 1.1703
40 2.7279 0.4738
41 1.6671 0.2847
42 0.5279 1.8008
43 2.5920 0.2097
44 0.1166 0.4072
45 1.2309 0.0818
46 1.5310 0.0495
47 2.3365 0.4658
48 2.7665 0.3874
49 1.4751 0.7557
50 3.0557 6.1322
51 4.8735 1.0582
52 1.9484 0.6160
53 3.8191 0.2127
54 2.5255 0.2452
55 2.0165 1.0315
56 2.1731 0.7328
57 2.6679 0.2711
58 1.6308 0.2126
59 3.3328 2.0916
60 18.4003 4.4689
61 1.7888 0.6419
62 2.7598 1.0068
63 2.8248 0.3748
64 2.3640 1.7492
65 1.8458 0.5081
66 0.6297 0.3991
67 5.3216 0.3283
68 3.1400 7.1256
69 0.5027 6.6879
70 2.6633 5.4521
71 1.1063 2.6034
72 3.5669 2.3752
73 0.5662 0.8380
74 9.2441 0.6662
75 1.5146 11.2440
76 4.2187 2.8334
77 6.0020 0.8937
78 1.2464 0.8364
79 1.5380 0.3340
80 2.6643 0.0425
81 4.7469 1.5871
82 1.1549 1.0175
83 3.5368 0.9786
84 1.5199 3.5657
85 1.2682 1.3209
86 1.8114 1.1837
87 2.5355 0.2638
88 5.9700 4.6171
89 2.1890 0.1371
90 3.3736 1.3964
91 2.1727 0.3579
92 2.3029 1.1651
93 1.7632 3.6888
94 2.2573 1.6589
95 2.1766 4.6732
96 1.8066 1.3340
97 1.2948 2.5728
98 3.6575 0.0340
99 1.9255 2.8553
100 2.6746 0.0409
101 4.0753 2.7106
102 3.7029 0.2403
103 1.2383 0.5040
104 3.3545 2.4060
105 1.9397 1.7027
106 2.4550 4.8171
107 2.0618 3.3052
108 1.3435 4.6498
109 2.2673 0.2051
110 3.5555 0.7126
111 1.2407 3.5422
112 1.5049 2.1250
113 3.3402 4.3515
114 2.8697 1.0321
115 2.7542 0.0040
116 7.2615 3.7939
117 3.2468 0.8441
118 1.8117 0.8938
119 4.3705 1.4611
120 2.1179 0.0955
121 1.7770 0.1557
122 20.0511 26.8629
123 6.6994 0.2037
124 0.3581 1.7189
125 2.5432 0.2537
126 1.9050 0.8637
127 0.9644 5.9414
128 8.8889 2.9256
129 2.7015 4.2652
130 1.6275 0.5036
131 1.3412 7.3181
132 1.6176 3.0569
133 1.8697 5.2481
134 1.1692 3.1464
135 2.1051 4.7705
136 2.8640 0.2179
137 3.2834 1.9727
138 5.4636 0.2856
139 2.7967 0.2343
140 2.2859 5.8578
141 1.9200 1.4742
142 1.8093 1.1257
143 2.5955 1.1074
144 1.3287 1.9300
145 2.6339 5.8340
146 2.0868 4.4240
147 5.5022 3.3783
148 3.3657 0.7014
149 2.8008 0.6823
150 3.6048 6.1489
151 3.4810 5.4528
152 3.1068 4.1442
153 2.9003 6.3911
154 4.0716 1.1922
155 2.6245 0.0731
156 4.0457 0.0008
157 3.6616 0.3105
158 4.8817 0.4513
159 6.0780 6.4378
160 3.0611 6.3942
161 2.3103 1.0533
162 1.6177 1.9111
163 2.2814 1.0962
164 2.1582 0.1997
165 3.4664 1.4434
166 3.0926 0.1488
167 4.0257 3.8463
168 3.4721 6.9908
169 3.1255 0.3502
170 5.2493 1.8816
171 2.6981 0.3604
172 4.3782 3.8207
173 3.7816 0.7036
174 4.4373 0.2978
175 2.6899 0.7978
176 5.1596 2.9269
177 3.5143 0.6850
178 5.9358 1.2195
}\tablaD

\title{\textbf{Taller -- Análisis de Componentes Principales (ACP)}\\[4pt]
\large Base \emph{Wine} (UCI Machine Learning Repository)}
\author{}
\date{\today}

\begin{document}
\maketitle

\noindent\textbf{Atribución del código.} El análisis reutiliza y adapta el código del cuaderno de clase:
Pacheco, M. (s.f.). \emph{Análisis Multivariado} [Repositorio de código y material de clase].
\url{https://datastatistic.github.io/multivariate/} (archivo \texttt{code/python/02-mrd/Código\_MRD\_1\_ACP.ipynb}).

\noindent\textbf{Datos.} Aeberhard, S. y Forina, M. (1992). \emph{Wine} [Dataset]. UCI Machine Learning Repository.
\url{https://doi.org/10.24432/C5PC7J}

\tableofcontents
\newpage

% =====================================================================
\section*{Justificación de la base}
\addcontentsline{toc}{section}{Justificación de la base}

\paragraph{(i) Contexto.} La base proviene del UCI Machine Learning Repository (donada en 1991; publicada por Forina et al., grupo PARVUS, Génova). Cada \textbf{fila es un vino} de la misma región de Italia, elaborado a partir de uno de \textbf{tres cultivares} (variedades de uva). Se midieron 13 constituyentes químicos por análisis de laboratorio. Hay $n = 178$ vinos, sin faltantes. Como la base no trae identificador, se crea \var{Vino\_001}, \dots, \var{Vino\_178}.

\begin{table}[H]
\centering\small
\caption{Diccionario de variables.}
\begin{tabular}{lll}
\toprule
Variable (notebook) & Original UCI & Descripción \\
\midrule
\var{alcohol} & Alcohol & Grado alcohólico (\% vol.) \\
\var{ac\_malico} & Malic acid & Ácido málico (acidez de origen frutal) \\
\var{ceniza} & Ash & Contenido de cenizas (minerales) \\
\var{alcal\_ceniza} & Alcalinity of ash & Alcalinidad de las cenizas \\
\var{magnesio} & Magnesium & Magnesio \\
\var{fenoles\_tot} & Total phenols & Fenoles totales \\
\var{flavonoides} & Flavanoids & Fenoles flavonoides \\
\var{fenoles\_noflav} & Nonflavanoid phenols & Fenoles no flavonoides \\
\var{proantocianinas} & Proanthocyanins & Proantocianidinas (taninos condensados) \\
\var{int\_color} & Color intensity & Intensidad de color \\
\var{tono} & Hue & Tonalidad del color \\
\var{od280\_od315} & OD280/OD315 & Razón de absorbancias (contenido proteico/fenólico) \\
\var{prolina} & Proline & Prolina (aminoácido) \\
\var{cultivar} & class & \textbf{Cualitativa}: cultivar 1, 2 o 3 \\
\bottomrule
\end{tabular}
\end{table}

\paragraph{(ii) Pregunta de investigación.} ¿Es posible resumir las 13 mediciones químicas de los vinos en unos pocos ejes interpretables, y esos ejes permiten diferenciar los tres cultivares de uva?

\paragraph{(iii) Pertinencia del ACP.} Varias variables miden aspectos del mismo fenómeno: fenoles totales, flavonoides, proantocianidinas y OD280/OD315 miden el contenido fenólico del vino, por lo que se espera que estén fuertemente correlacionadas (redundancia). Alcohol, color y prolina también se asocian con la madurez de la uva. Esa redundancia es la que el ACP puede resumir.

% =====================================================================
\section{Carga y exploración inicial de los datos}

La base tiene 178 vinos y 14 columnas: 13 variables cuantitativas continuas y una cualitativa (\var{cultivar}). No hay valores faltantes en ninguna variable (0\,\%).

\paragraph{Reflexión.} Cada fila de la base representa el análisis de un vino específico y cada columna una característica química cuantitativa continua. Las variables cuantitativas se consideran \emph{activas} porque intervienen directamente en el cálculo de las componentes principales y de la varianza total, mientras que el cultivar se considera \emph{suplementaria}: no participa en la construcción de los ejes, pero se usa para interpretarlos.

% =====================================================================
\section{Análisis descriptivo univariado y matriz de correlaciones}

\begin{table}[H]
\centering\small
\caption{Estadísticos descriptivos, coeficiente de variación (CV) y asimetría.}
\begin{tabular}{lrrrrrrr}
\toprule
Variable & Media & Desv. est. & Mín. & Mediana & Máx. & CV (\%) & Asimetría \\
\midrule
\texttt{alcohol} & 13.001 & 0.812 & 11.03 & 13.050 & 14.83 & 6.2 & -0.05 \\
\texttt{ac\_malico} & 2.336 & 1.117 & 0.74 & 1.865 & 5.80 & 47.8 & 1.03 \\
\texttt{ceniza} & 2.367 & 0.274 & 1.36 & 2.360 & 3.23 & 11.6 & -0.18 \\
\texttt{alcal\_ceniza} & 19.495 & 3.340 & 10.60 & 19.500 & 30.00 & 17.1 & 0.21 \\
\texttt{magnesio} & 99.742 & 14.282 & 70.00 & 98.000 & 162.00 & 14.3 & 1.09 \\
\texttt{fenoles\_tot} & 2.295 & 0.626 & 0.98 & 2.355 & 3.88 & 27.3 & 0.09 \\
\texttt{flavonoides} & 2.029 & 0.999 & 0.34 & 2.135 & 5.08 & 49.2 & 0.03 \\
\texttt{fenoles\_noflav} & 0.362 & 0.124 & 0.13 & 0.340 & 0.66 & 34.4 & 0.45 \\
\texttt{proantocianinas} & 1.591 & 0.572 & 0.41 & 1.555 & 3.58 & 36.0 & 0.51 \\
\texttt{int\_color} & 5.058 & 2.318 & 1.28 & 4.690 & 13.00 & 45.8 & 0.86 \\
\texttt{tono} & 0.957 & 0.229 & 0.48 & 0.965 & 1.71 & 23.9 & 0.02 \\
\texttt{od280\_od315} & 2.612 & 0.710 & 1.27 & 2.780 & 4.00 & 27.2 & -0.30 \\
\texttt{prolina} & 746.893 & 314.907 & 278.00 & 673.500 & 1680.00 & 42.2 & 0.76 \\
\bottomrule
\end{tabular}
\end{table}

Las escalas son muy distintas: \var{prolina} tiene media $\approx 747$ y desviación $\approx 315$, \var{magnesio} $\approx 100$, mientras que \var{fenoles\_noflav} o \var{tono} están por debajo de 1. Los mayores coeficientes de variación los tienen \var{flavonoides} ($\approx 49\,\%$), \var{ac\_malico} ($\approx 48\,\%$) e \var{int\_color} ($\approx 46\,\%$); los menores, \var{alcohol} ($\approx 6\,\%$) y \var{ceniza} ($\approx 12\,\%$).

\begin{figure}[H]
\centering
\begin{tikzpicture}
\begin{groupplot}[
  group style={group size=4 by 4, horizontal sep=1.0cm, vertical sep=1.25cm},
  width=4.6cm, height=3.6cm, ejes, ytick=\empty, xticklabel style={/pgf/number format/fixed}, max space between ticks=22,
  tick label style={font=\tiny}, title style={font=\scriptsize\ttfamily}]
\nextgroupplot[title=alcohol]        \addplot[ybar interval, hist={bins=14}, fill=azul!60, draw=azul] table[y=alcohol]{\datos};
\nextgroupplot[title=ac\_malico]     \addplot[ybar interval, hist={bins=14}, fill=azul!60, draw=azul] table[y=acmalico]{\datos};
\nextgroupplot[title=ceniza]         \addplot[ybar interval, hist={bins=14}, fill=azul!60, draw=azul] table[y=ceniza]{\datos};
\nextgroupplot[title=alcal\_ceniza]  \addplot[ybar interval, hist={bins=14}, fill=azul!60, draw=azul] table[y=alcalceniza]{\datos};
\nextgroupplot[title=magnesio]       \addplot[ybar interval, hist={bins=14}, fill=azul!60, draw=azul] table[y=magnesio]{\datos};
\nextgroupplot[title=fenoles\_tot]   \addplot[ybar interval, hist={bins=14}, fill=azul!60, draw=azul] table[y=fenolestot]{\datos};
\nextgroupplot[title=flavonoides]    \addplot[ybar interval, hist={bins=14}, fill=azul!60, draw=azul] table[y=flavonoides]{\datos};
\nextgroupplot[title=fenoles\_noflav]\addplot[ybar interval, hist={bins=14}, fill=azul!60, draw=azul] table[y=fenolesnoflav]{\datos};
\nextgroupplot[title=proantocianinas]\addplot[ybar interval, hist={bins=14}, fill=azul!60, draw=azul] table[y=proanto]{\datos};
\nextgroupplot[title=int\_color]     \addplot[ybar interval, hist={bins=14}, fill=azul!60, draw=azul] table[y=intcolor]{\datos};
\nextgroupplot[title=tono]           \addplot[ybar interval, hist={bins=14}, fill=azul!60, draw=azul] table[y=tono]{\datos};
\nextgroupplot[title=od280\_od315]   \addplot[ybar interval, hist={bins=14}, fill=azul!60, draw=azul] table[y=odratio]{\datos};
\nextgroupplot[title=prolina]        \addplot[ybar interval, hist={bins=14}, fill=azul!60, draw=azul] table[y=prolina]{\datos};
\end{groupplot}
\end{tikzpicture}
\caption{Histogramas de las 13 variables químicas.}
\label{fig:hist}
\end{figure}

\begin{figure}[H]
\centering
\begin{tikzpicture}
\begin{axis}[ejes, width=\textwidth, height=7cm,
  boxplot/draw direction=y, boxplot/box extend=0.55,
  xtick={1,...,13}, xticklabels={\texttt{alcohol},\texttt{ac\_malico},\texttt{ceniza},\texttt{alcal\_ceniza},\texttt{magnesio},\texttt{fenoles\_tot},\texttt{flavonoides},\texttt{fenoles\_noflav},\texttt{proantocianinas},\texttt{int\_color},\texttt{tono},\texttt{od280\_od315},\texttt{prolina}},
  x tick label style={rotate=55, anchor=north east, font=\scriptsize},
  ylabel={Valor estandarizado}, xmin=0.3, xmax=13.7,
  title={Diagramas de caja (variables estandarizadas)}]
\addplot+[boxplot prepared={median=0.061,lower quartile=-0.786,upper quartile=0.834,lower whisker=-2.427,upper whisker=2.253}, fill=caja, draw=black, mark=*, mark options={fill=rojo,draw=rojo,scale=.6}] coordinates {};
\addplot+[boxplot prepared={median=-0.422,lower quartile=-0.657,upper quartile=0.668,lower whisker=-1.429,upper whisker=2.554}, fill=caja, draw=black, mark=*, mark options={fill=rojo,draw=rojo,scale=.6}] coordinates {(0,3.100) (0,2.841) (0,2.966)};
\addplot+[boxplot prepared={median=-0.024,lower quartile=-0.571,upper quartile=0.696,lower whisker=-2.429,upper whisker=2.017}, fill=caja, draw=black, mark=*, mark options={fill=rojo,draw=rojo,scale=.6}] coordinates {(0,3.111) (0,-3.669) (0,3.147)};
\addplot+[boxplot prepared={median=0.002,lower quartile=-0.687,upper quartile=0.600,lower whisker=-2.484,upper whisker=2.247}, fill=caja, draw=black, mark=*, mark options={fill=rojo,draw=rojo,scale=.6}] coordinates {(0,-2.664) (0,3.146) (0,2.696) (0,2.696)};
\addplot+[boxplot prepared={median=-0.122,lower quartile=-0.822,upper quartile=0.508,lower whisker=-2.082,upper whisker=2.399}, fill=caja, draw=black, mark=*, mark options={fill=rojo,draw=rojo,scale=.6}] coordinates {(0,3.589) (0,2.749) (0,2.539) (0,4.359)};
\addplot+[boxplot prepared={median=0.096,lower quartile=-0.883,upper quartile=0.807,lower whisker=-2.101,upper whisker=2.532}, fill=caja, draw=black, mark=*, mark options={fill=rojo,draw=rojo,scale=.6}] coordinates {};
\addplot+[boxplot prepared={median=0.106,lower quartile=-0.825,upper quartile=0.847,lower whisker=-1.691,upper whisker=3.054}, fill=caja, draw=black, mark=*, mark options={fill=rojo,draw=rojo,scale=.6}] coordinates {};
\addplot+[boxplot prepared={median=-0.176,lower quartile=-0.738,upper quartile=0.608,lower whisker=-1.863,upper whisker=2.396}, fill=caja, draw=black, mark=*, mark options={fill=rojo,draw=rojo,scale=.6}] coordinates {};
\addplot+[boxplot prepared={median=-0.063,lower quartile=-0.596,upper quartile=0.627,lower whisker=-2.063,upper whisker=2.392}, fill=caja, draw=black, mark=*, mark options={fill=rojo,draw=rojo,scale=.6}] coordinates {(0,2.951) (0,3.475)};
\addplot+[boxplot prepared={median=-0.159,lower quartile=-0.793,upper quartile=0.493,lower whisker=-1.630,upper whisker=2.356}, fill=caja, draw=black, mark=*, mark options={fill=rojo,draw=rojo,scale=.6}] coordinates {(0,2.477) (0,3.426) (0,2.887) (0,2.425)};
\addplot+[boxplot prepared={median=0.033,lower quartile=-0.765,upper quartile=0.711,lower whisker=-2.089,upper whisker=2.155}, fill=caja, draw=black, mark=*, mark options={fill=rojo,draw=rojo,scale=.6}] coordinates {(0,3.292)};
\addplot+[boxplot prepared={median=0.237,lower quartile=-0.950,upper quartile=0.786,lower whisker=-1.890,upper whisker=1.955}, fill=caja, draw=black, mark=*, mark options={fill=rojo,draw=rojo,scale=.6}] coordinates {};
\addplot+[boxplot prepared={median=-0.233,lower quartile=-0.782,upper quartile=0.756,lower whisker=-1.489,upper whisker=2.963}, fill=caja, draw=black, mark=*, mark options={fill=rojo,draw=rojo,scale=.6}] coordinates {};
\addplot[cero] coordinates {(0.3,0) (13.7,0)};
\end{axis}
\end{tikzpicture}
\caption{Diagramas de caja de las variables estandarizadas; los puntos rojos son atípicos según la regla de $1.5\cdot\mathrm{RIC}$.}
\label{fig:box}
\end{figure}

\var{ac\_malico} (asimetría $\approx 1.0$), \var{magnesio} ($\approx 1.1$), \var{int\_color} ($\approx 0.9$) y \var{prolina} ($\approx 0.8$) tienen asimetría positiva (cola a la derecha); el resto es aproximadamente simétrico ($|\text{asimetría}| < 0.5$). Por la regla de $1.5\cdot\mathrm{RIC}$ hay pocos atípicos univariados: 4 en \var{alcal\_ceniza}, \var{magnesio} e \var{int\_color}; 3 en \var{ac\_malico} y \var{ceniza}; 2 en \var{proantocianinas} y 1 en \var{tono}. Los valores extremos son químicamente plausibles (no parecen errores de digitación), por lo que se conservan.

\begin{figure}[H]
\centering
\begin{tikzpicture}
\begin{axis}[
  width=13cm, height=13cm, axis on top, enlargelimits=false,
  colormap={rdbu}{rgb255=(33,102,172) rgb255=(146,197,222) rgb255=(247,247,247) rgb255=(244,165,130) rgb255=(178,24,43)},
  point meta min=-1, point meta max=1,
  colorbar, colorbar style={ylabel={$r$}, ytick={-1,-0.5,0,0.5,1}},
  xtick={0,...,12}, ytick={0,...,12},
  xticklabels={\texttt{alcohol},\texttt{ac\_malico},\texttt{ceniza},\texttt{alcal\_ceniza},\texttt{magnesio},\texttt{fenoles\_tot},\texttt{flavonoides},\texttt{fenoles\_noflav},\texttt{proantocianinas},\texttt{int\_color},\texttt{tono},\texttt{od280\_od315},\texttt{prolina}}, yticklabels={\texttt{prolina},\texttt{od280\_od315},\texttt{tono},\texttt{int\_color},\texttt{proantocianinas},\texttt{fenoles\_noflav},\texttt{flavonoides},\texttt{fenoles\_tot},\texttt{magnesio},\texttt{alcal\_ceniza},\texttt{ceniza},\texttt{ac\_malico},\texttt{alcohol}},
  x tick label style={rotate=60, anchor=north east, font=\scriptsize},
  y tick label style={font=\scriptsize},
  title={Matriz de correlaciones $R_x$},
  nodes near coords={\pgfmathprintnumber[fixed,precision=2,fixed zerofill]\pgfplotspointmeta},
  every node near coord/.append style={font=\tiny, anchor=center, black},
]
\addplot[matrix plot*, mesh/cols=13, point meta=explicit] table[meta=c]{
x y c
0 12 1.00
1 12 0.09
2 12 0.21
3 12 -0.31
4 12 0.27
5 12 0.29
6 12 0.24
7 12 -0.16
8 12 0.14
9 12 0.55
10 12 -0.07
11 12 0.07
12 12 0.64
0 11 0.09
1 11 1.00
2 11 0.16
3 11 0.29
4 11 -0.05
5 11 -0.34
6 11 -0.41
7 11 0.29
8 11 -0.22
9 11 0.25
10 11 -0.56
11 11 -0.37
12 11 -0.19
0 10 0.21
1 10 0.16
2 10 1.00
3 10 0.44
4 10 0.29
5 10 0.13
6 10 0.12
7 10 0.19
8 10 0.01
9 10 0.26
10 10 -0.07
11 10 0.00
12 10 0.22
0 9 -0.31
1 9 0.29
2 9 0.44
3 9 1.00
4 9 -0.08
5 9 -0.32
6 9 -0.35
7 9 0.36
8 9 -0.20
9 9 0.02
10 9 -0.27
11 9 -0.28
12 9 -0.44
0 8 0.27
1 8 -0.05
2 8 0.29
3 8 -0.08
4 8 1.00
5 8 0.21
6 8 0.20
7 8 -0.26
8 8 0.24
9 8 0.20
10 8 0.06
11 8 0.07
12 8 0.39
0 7 0.29
1 7 -0.34
2 7 0.13
3 7 -0.32
4 7 0.21
5 7 1.00
6 7 0.86
7 7 -0.45
8 7 0.61
9 7 -0.06
10 7 0.43
11 7 0.70
12 7 0.50
0 6 0.24
1 6 -0.41
2 6 0.12
3 6 -0.35
4 6 0.20
5 6 0.86
6 6 1.00
7 6 -0.54
8 6 0.65
9 6 -0.17
10 6 0.54
11 6 0.79
12 6 0.49
0 5 -0.16
1 5 0.29
2 5 0.19
3 5 0.36
4 5 -0.26
5 5 -0.45
6 5 -0.54
7 5 1.00
8 5 -0.37
9 5 0.14
10 5 -0.26
11 5 -0.50
12 5 -0.31
0 4 0.14
1 4 -0.22
2 4 0.01
3 4 -0.20
4 4 0.24
5 4 0.61
6 4 0.65
7 4 -0.37
8 4 1.00
9 4 -0.03
10 4 0.30
11 4 0.52
12 4 0.33
0 3 0.55
1 3 0.25
2 3 0.26
3 3 0.02
4 3 0.20
5 3 -0.06
6 3 -0.17
7 3 0.14
8 3 -0.03
9 3 1.00
10 3 -0.52
11 3 -0.43
12 3 0.32
0 2 -0.07
1 2 -0.56
2 2 -0.07
3 2 -0.27
4 2 0.06
5 2 0.43
6 2 0.54
7 2 -0.26
8 2 0.30
9 2 -0.52
10 2 1.00
11 2 0.57
12 2 0.24
0 1 0.07
1 1 -0.37
2 1 0.00
3 1 -0.28
4 1 0.07
5 1 0.70
6 1 0.79
7 1 -0.50
8 1 0.52
9 1 -0.43
10 1 0.57
11 1 1.00
12 1 0.31
0 0 0.64
1 0 -0.19
2 0 0.22
3 0 -0.44
4 0 0.39
5 0 0.50
6 0 0.49
7 0 -0.31
8 0 0.33
9 0 0.32
10 0 0.24
11 0 0.31
12 0 1.00
};
\end{axis}
\end{tikzpicture}
\caption{Mapa de calor de la matriz de correlaciones $R_x$.}
\label{fig:heat}
\end{figure}

\paragraph{Interpretación por bloques.}
\begin{itemize}[leftmargin=*]
  \item \textbf{Bloque fenólico:} \var{fenoles\_tot}, \var{flavonoides}, \var{proantocianinas} y \var{od280\_od315} tienen correlaciones positivas altas entre sí ($\approx 0.6$--$0.86$; flavonoides--fenoles totales $\approx 0.86$). \var{tono} también se une a este bloque; \var{fenoles\_noflav} y \var{ac\_malico} se correlacionan negativamente con él.
  \item \textbf{Bloque alcohol/color/prolina:} \var{alcohol}, \var{int\_color} y \var{prolina} se correlacionan positivamente (alcohol--prolina $\approx 0.64$; alcohol--color $\approx 0.55$).
  \item \textbf{Bloque mineral:} \var{ceniza} y \var{alcal\_ceniza} se correlacionan moderadamente ($\approx 0.44$) y casi no se relacionan con el resto.
\end{itemize}

\paragraph{Reflexión.} Al observar $R_x$ se identifican aproximadamente 3 o 4 factores principales que agrupan y explican la variabilidad de los atributos químicos del vino.

% =====================================================================
\section{Preprocesamiento: faltantes y estandarización}

Con \texttt{dropna()} se pierden \textbf{0 observaciones} (la base no tiene faltantes) y se conservan $178 \geq 100$. Se construye la matriz estandarizada y ponderada
\[
  z_{ij} = \frac{x_{ij} - \bar{x}_j}{s_j\sqrt{n-1}}, \qquad i = 1,\dots,178,\; j = 1,\dots,13,
\]
que cumple $Z^\top Z = R_x$ (diferencia máxima $\approx 1.6\times 10^{-15}$, es decir, error de redondeo).

\begin{table}[H]
\centering\small
\caption{Varianza de cada variable sin estandarizar y su peso en la traza de la matriz de covarianzas.}
\begin{tabular}{lrr}
\toprule
Variable & Varianza & \% de la varianza total \\
\midrule
\var{prolina} & 99\,166.717 & 99.774 \\
\var{magnesio} & 203.989 & 0.205 \\
\var{alcal\_ceniza} & 11.153 & 0.011 \\
\var{int\_color} & 5.374 & 0.005 \\
Resto (9 variables) & $< 1.3$ & $< 0.005$ \\
\bottomrule
\end{tabular}
\end{table}

\paragraph{Reflexión.} Es necesario estandarizar porque las variables están en unidades y órdenes de magnitud muy dispares. Sin estandarizar, la matriz de covarianzas estaría dominada por \var{prolina} (99.8\,\% de la varianza total) y \var{magnesio}; el ACP reflejaría casi únicamente la variabilidad de esas variables y no permitiría interpretar las demás.

% =====================================================================
\section{Verificación de supuestos y pertinencia del ACP}

\begin{figure}[H]
\centering
\begin{tikzpicture}
\begin{groupplot}[group style={group size=3 by 1, horizontal sep=1.4cm},
  width=5.6cm, height=5.2cm, ejes, every axis plot/.append style={only marks, mark size=1.3pt}]
\nextgroupplot[xlabel=\var{fenoles\_tot}, ylabel=\var{flavonoides}, legend to name=leyCult, legend columns=3,
  legend style={font=\footnotesize, draw=none}]
\addplot[c1, mark=*, opacity=.75, discard if not={cult}{1}] table[x=fenolestot, y=flavonoides]{\datos};
\addplot[c2, mark=*, opacity=.75, discard if not={cult}{2}] table[x=fenolestot, y=flavonoides]{\datos};
\addplot[c3, mark=*, opacity=.75, discard if not={cult}{3}] table[x=fenolestot, y=flavonoides]{\datos};
\legend{Cultivar 1, Cultivar 2, Cultivar 3}
\nextgroupplot[xlabel=\var{alcohol}, ylabel=\var{prolina}]
\addplot[c1, mark=*, opacity=.75, discard if not={cult}{1}] table[x=alcohol, y=prolina]{\datos};
\addplot[c2, mark=*, opacity=.75, discard if not={cult}{2}] table[x=alcohol, y=prolina]{\datos};
\addplot[c3, mark=*, opacity=.75, discard if not={cult}{3}] table[x=alcohol, y=prolina]{\datos};
\nextgroupplot[xlabel=\var{ac\_malico}, ylabel=\var{tono}]
\addplot[c1, mark=*, opacity=.75, discard if not={cult}{1}] table[x=acmalico, y=tono]{\datos};
\addplot[c2, mark=*, opacity=.75, discard if not={cult}{2}] table[x=acmalico, y=tono]{\datos};
\addplot[c3, mark=*, opacity=.75, discard if not={cult}{3}] table[x=acmalico, y=tono]{\datos};
\end{groupplot}
\end{tikzpicture}

\ref*{leyCult}
\caption{Diagramas de dispersión de pares representativos (bloque fenólico y variables asimétricas).}
\label{fig:pairs}
\end{figure}

Las relaciones dentro del bloque fenólico son claramente lineales y no se observan relaciones curvas marcadas. Las variables asimétricas (\var{prolina}, \var{int\_color}, \var{ac\_malico}, \var{magnesio}) muestran nubes en abanico (más dispersión en valores altos). La Tabla~\ref{tab:log} lo comprueba: al aplicar logaritmo la correlación de Pearson casi no cambia y coincide con la de Spearman, es decir, las relaciones son monótonas y prácticamente lineales. Se trabaja en la escala original.

\begin{table}[H]
\centering\small
\caption{¿Mejora la linealidad con logaritmo? Pearson vs.\ Spearman.}
\label{tab:log}
\begin{tabular}{lrrr}
\toprule
Par & Pearson original & Pearson (log) & Spearman \\
\midrule
\var{prolina}--\var{alcohol} & 0.644 & 0.637 & 0.634 \\
\var{ac\_malico}--\var{tono} & $-0.561$ & $-0.562$ & $-0.560$ \\
\bottomrule
\end{tabular}
\end{table}

\paragraph{Determinante de $R_x$.} $\det(R_x) = 0.000469 \approx 0.00047$, muy cercano a 0 (el máximo posible, con variables incorrelacionadas, sería 1). La varianza generalizada es muy pequeña: la nube de puntos en 13 dimensiones está ``aplastada'' en algunas direcciones porque hay mucha redundancia lineal entre las variables. Esto justifica el ACP.

% =====================================================================
\section{Cálculo de las componentes principales ``a mano'' (DVS)}

Se aplica la descomposición en valores singulares
\[
  Z = P\,\Delta\,Q^\top, \qquad \lambda_k = \delta_k^2, \qquad F = Z\,Q = P\,\Delta,
\]
donde $Q$ contiene los vectores propios de $Z^\top Z = R_x$ y $P$ los de $Z Z^\top$. Se verificaron numéricamente todas las propiedades:
\begin{itemize}[leftmargin=*]
  \item $F = P\Delta$ y $Q^\top Q = I_{13}$.
  \item $\sum_k \lambda_k = 13 = p$.
  \item $\operatorname{Var}(F_k) = \lambda_k/(n-1)$, es decir, proporcional a $\lambda_k$ con constante $1/177$.
  \item Las componentes son incorrelacionadas: $\max_{k\neq l}|\operatorname{cor}(F_k,F_l)| \approx 1.4\times10^{-15}$.
\end{itemize}

Los valores propios de la DVS coinciden con los de \texttt{eigh(cor(X))}. Los vectores coinciden en valor absoluto, pero algunas columnas salen con \textbf{signo cambiado}: si $q_k$ es vector propio de $R_x$, también lo es $-q_k$, pues $R_x(-q_k) = \lambda_k(-q_k)$ y $\|-q_k\| = 1$. Cada algoritmo elige el signo de forma arbitraria; invertirlo solo refleja el eje y no cambia la varianza ni la interpretación, siempre que se lean juntos cargas y coordenadas.

% =====================================================================
\section{Selección del número de componentes}

\begin{table}[H]
\centering\small
\caption{Valores propios e inercia.}
\begin{tabular}{lrrr}
\toprule
Componente & $\lambda_k$ & Inercia (\%) & Inercia acumulada (\%) \\
\midrule
CP1 & 4.706 & 36.199 & 36.199 \\
CP2 & 2.497 & 19.207 & 55.406 \\
CP3 & 1.446 & 11.124 & 66.530 \\
CP4 & 0.919 & 7.069 & 73.599 \\
CP5 & 0.853 & 6.563 & 80.162 \\
CP6 & 0.642 & 4.936 & 85.098 \\
CP7 & 0.551 & 4.239 & 89.337 \\
CP8 & 0.348 & 2.681 & 92.018 \\
CP9 & 0.289 & 2.222 & 94.240 \\
CP10 & 0.251 & 1.930 & 96.170 \\
CP11 & 0.226 & 1.737 & 97.907 \\
CP12 & 0.169 & 1.298 & 99.205 \\
CP13 & 0.103 & 0.795 & 100.000 \\
\bottomrule
\end{tabular}
\end{table}

\begin{figure}[H]
\centering
\begin{tikzpicture}
\begin{axis}[ejes, width=10cm, height=6.5cm, ybar, bar width=14pt,
  xtick={1,...,13}, xlabel={Componente}, ylabel={\% de inercia},
  ymin=0, ymax=42, xmin=0.3, xmax=13.7,
  nodes near coords={\pgfmathprintnumber[fixed,precision=1]\pgfplotspointmeta},
  every node near coord/.append style={font=\tiny, yshift=5pt},
  legend style={font=\footnotesize, at={(0.97,0.97)}, anchor=north east},
  title={Scree plot -- criterio del codo}]
\addplot[fill=azul!70, draw=azul] coordinates {(1,36.20) (2,19.21) (3,11.12) (4,7.07) (5,6.56) (6,4.94) (7,4.24) (8,2.68) (9,2.22) (10,1.93) (11,1.74) (12,1.30) (13,0.80)};
\addplot[sharp plot, black, mark=*, mark size=1.5pt, nodes near coords={}] coordinates {(1,36.20) (2,19.21) (3,11.12) (4,7.07) (5,6.56) (6,4.94) (7,4.24) (8,2.68) (9,2.22) (10,1.93) (11,1.74) (12,1.30) (13,0.80)};
\addplot[sharp plot, rojo, dashed, thick, nodes near coords={}] coordinates {(0.3,7.692) (13.7,7.692)};
\legend{Inercia, , Kaiser ($\lambda=1 \leftrightarrow 100/p\,\%$)}
\end{axis}
\end{tikzpicture}
\hfill
\begin{tikzpicture}
\begin{axis}[ejes, width=6.2cm, height=6.5cm, xlabel={Número de componentes}, ylabel={GCV},
  xtick={1,...,8}, title={Validación cruzada (GCV)}]
\addplot[azul, thick, mark=*] coordinates {(1,0.7616) (2,0.6408) (3,0.5886) (4,0.5799) (5,0.5579) (6,0.5538) (7,0.5457) (8,0.5953)};
\end{axis}
\end{tikzpicture}
\caption{Izquierda: porcentaje de inercia por componente con la línea de Kaiser. Derecha: criterio GCV.}
\label{fig:scree}
\end{figure}

\begin{table}[H]
\centering\small
\caption{Comparación de criterios.}
\begin{tabular}{lcc}
\toprule
Criterio & $k$ sugerido & Inercia acumulada con ese $k$ (\%) \\
\midrule
Kaiser ($\lambda > 1$) & 3 & 66.5 \\
Inercia acumulada $\geq 70\,\%$ & 4 & 73.6 \\
Codo (scree plot) & 3 & 66.5 \\
GCV (validación cruzada) & 7 & 89.3 \\
\bottomrule
\end{tabular}
\end{table}

\paragraph{Decisión.} Kaiser y el codo sugieren \textbf{$k = 3$}, que explica $\approx 66.5\,\%$ de la inercia, apenas por debajo del 70\,\%. Llegar al 70\,\% exige una cuarta componente con $\lambda_4 \approx 0.92 < 1$: aporta menos que una sola variable original y su inercia ($\approx 7.1\,\%$) es casi igual a la de CP5 ($\approx 6.6\,\%$), por lo que no forma un eje estable. El GCV sugiere 7, pero su curva es casi plana desde $k=3$ ($0.589 \to 0.546$) y CP4--CP7 tienen $\lambda < 1$. Se retienen \textbf{$k = 3$ componentes} (CP1 $\approx 36.2\,\%$, CP2 $\approx 19.2\,\%$, CP3 $\approx 11.1\,\%$).

\paragraph{Reflexión.} Cuando los criterios no coinciden se debe priorizar el porcentaje de varianza acumulada junto con la interpretabilidad. Si un número de componentes retiene un porcentaje satisfactorio de la inercia y permite una interpretación química clara respecto a la pregunta de investigación, es preferible, aunque un criterio estricto sugiera añadir una dimensión más.

% =====================================================================
\section{Visualización}

\begin{figure}[H]
\centering
\begin{tikzpicture}
\begin{axis}[ejes, width=10.5cm, height=8.5cm,
  xlabel={CP1 (36.2\,\%)}, ylabel={CP2 (19.2\,\%)},
  legend style={font=\footnotesize, at={(0.98,0.02)}, anchor=south east},
  title={Primer plano principal de los vinos ($F = ZQ$)}]
\addplot[only marks, mark=*, mark size=1.6pt, c1, discard if not={cult}{1}] table[x=cpa, y=cpb]{\tablaF};
\addplot[only marks, mark=*, mark size=1.6pt, c2, discard if not={cult}{2}] table[x=cpa, y=cpb]{\tablaF};
\addplot[only marks, mark=*, mark size=1.6pt, c3, discard if not={cult}{3}] table[x=cpa, y=cpb]{\tablaF};
\legend{Cultivar 1, Cultivar 2, Cultivar 3}
\draw[gray, dashed] (axis cs:-0.4,0) -- (axis cs:0.4,0);
\draw[gray, dashed] (axis cs:0,-0.4) -- (axis cs:0,0.4);
\node[font=\tiny\ttfamily, anchor=south west, inner sep=1pt] at (axis cs:-0.0361,0.2902) {Vino\_116};
\node[font=\tiny\ttfamily, anchor=south west, inner sep=1pt] at (axis cs:-0.2816,-0.2066) {Vino\_004};
\node[font=\tiny\ttfamily, anchor=south west, inner sep=1pt] at (axis cs:-0.3233,-0.1571) {Vino\_015};
\node[font=\tiny\ttfamily, anchor=south west, inner sep=1pt] at (axis cs:0.2405,-0.2075) {Vino\_178};
\node[font=\tiny\ttfamily, anchor=south west, inner sep=1pt] at (axis cs:-0.2655,-0.1888) {Vino\_019};
\node[font=\tiny\ttfamily, anchor=south west, inner sep=1pt] at (axis cs:0.0786,-0.2635) {Vino\_159};
\node[font=\tiny\ttfamily, anchor=south west, inner sep=1pt] at (axis cs:-0.0571,0.2530) {Vino\_081};
\node[font=\tiny\ttfamily, anchor=south west, inner sep=1pt] at (axis cs:0.2007,-0.2069) {Vino\_176};
\end{axis}
\end{tikzpicture}
\caption{Proyección de los vinos en el primer plano principal; se etiquetan los 8 vinos con mayor contribución.}
\label{fig:ind12}
\end{figure}

\begin{figure}[H]
\centering
\begin{tikzpicture}
\begin{axis}[width=11cm, height=11cm, ejes, axis equal image,
  xmin=-1.25, xmax=1.25, ymin=-1.25, ymax=1.25, xtick={-1,-0.5,0,0.5,1}, ytick={-1,-0.5,0,0.5,1},
  xlabel={CP1 (36.2\,\%)}, ylabel={CP2 (19.2\,\%)}, title={Plano (1,2)}]
\draw[black] (axis cs:0,0) circle[radius=1];
\draw[gray, dashed] (axis cs:-1.25,0) -- (axis cs:1.25,0); \draw[gray, dashed] (axis cs:0,-1.25) -- (axis cs:0,1.25);
\draw[-{Stealth[length=2mm]}, thick, color=azulcos!68!gray] (axis cs:0,0) -- (axis cs:-0.313,-0.764) node[anchor=68, font=\scriptsize, inner sep=1pt] {\texttt{alcohol}};
\draw[-{Stealth[length=2mm]}, thick, color=azulcos!41!gray] (axis cs:0,0) -- (axis cs:0.532,-0.355) node[anchor=146, font=\scriptsize, inner sep=1pt] {\texttt{ac\_malico}};
\draw[-{Stealth[length=2mm]}, thick, color=azulcos!25!gray] (axis cs:0,0) -- (axis cs:0.004,-0.499) node[anchor=91, font=\scriptsize, inner sep=1pt] {\texttt{ceniza}};
\draw[-{Stealth[length=2mm]}, thick, color=azulcos!27!gray] (axis cs:0,0) -- (axis cs:0.519,0.017) node[anchor=182, font=\scriptsize, inner sep=1pt] {\texttt{alcal\_ceniza}};
\draw[-{Stealth[length=2mm]}, thick, color=azulcos!32!gray] (axis cs:0,0) -- (axis cs:-0.308,-0.473) node[anchor=57, font=\scriptsize, inner sep=1pt] {\texttt{magnesio}};
\draw[-{Stealth[length=2mm]}, thick, color=azulcos!74!gray] (axis cs:0,0) -- (axis cs:-0.856,-0.103) node[anchor=7, font=\scriptsize, inner sep=1pt] {\texttt{fenoles\_tot}};
\draw[-{Stealth[length=2mm]}, thick, color=azulcos!84!gray] (axis cs:0,0) -- (axis cs:-0.917,0.005) node[anchor=360, font=\scriptsize, inner sep=1pt] {\texttt{flavonoides}};
\draw[-{Stealth[length=2mm]}, thick, color=azulcos!42!gray] (axis cs:0,0) -- (axis cs:0.648,-0.045) node[anchor=176, font=\scriptsize, inner sep=1pt] {\texttt{fenoles\_noflav}};
\draw[-{Stealth[length=2mm]}, thick, color=azulcos!47!gray] (axis cs:0,0) -- (axis cs:-0.680,-0.062) node[anchor=5, font=\scriptsize, inner sep=1pt] {\texttt{proantocianinas}};
\draw[-{Stealth[length=2mm]}, thick, color=azulcos!74!gray] (axis cs:0,0) -- (axis cs:0.192,-0.837) node[anchor=103, font=\scriptsize, inner sep=1pt] {\texttt{int\_color}};
\draw[-{Stealth[length=2mm]}, thick, color=azulcos!61!gray] (axis cs:0,0) -- (axis cs:-0.644,0.441) node[anchor=326, font=\scriptsize, inner sep=1pt] {\texttt{tono}};
\draw[-{Stealth[length=2mm]}, thick, color=azulcos!73!gray] (axis cs:0,0) -- (axis cs:-0.816,0.260) node[anchor=342, font=\scriptsize, inner sep=1pt] {\texttt{od280\_od315}};
\draw[-{Stealth[length=2mm]}, thick, color=azulcos!72!gray] (axis cs:0,0) -- (axis cs:-0.622,-0.577) node[anchor=43, font=\scriptsize, inner sep=1pt] {\texttt{prolina}};
\end{axis}
\end{tikzpicture}

\vspace{4pt}
\begin{tikzpicture}
\begin{groupplot}[group style={group size=2 by 1, horizontal sep=1.8cm},
  width=8cm, height=8cm, ejes, axis equal image,
  xmin=-1.25, xmax=1.25, ymin=-1.25, ymax=1.25, xtick={-1,0,1}, ytick={-1,0,1}]
\nextgroupplot[xlabel={CP1 (36.2\,\%)}, ylabel={CP3 (11.1\,\%)}, title={Plano (1,3)}]
\draw[black] (axis cs:0,0) circle[radius=1];
\draw[gray, dashed] (axis cs:-1.25,0) -- (axis cs:1.25,0); \draw[gray, dashed] (axis cs:0,-1.25) -- (axis cs:0,1.25);
\draw[-{Stealth[length=2mm]}, thick, color=azulcos!16!gray] (axis cs:0,0) -- (axis cs:-0.313,-0.249) node[anchor=39, font=\scriptsize, inner sep=1pt] {\texttt{alcohol}};
\draw[-{Stealth[length=2mm]}, thick, color=azulcos!29!gray] (axis cs:0,0) -- (axis cs:0.532,0.107) node[anchor=191, font=\scriptsize, inner sep=1pt] {\texttt{ac\_malico}};
\draw[-{Stealth[length=2mm]}, thick, color=azulcos!57!gray] (axis cs:0,0) -- (axis cs:0.004,0.753) node[anchor=270, font=\scriptsize, inner sep=1pt] {\texttt{ceniza}};
\draw[-{Stealth[length=2mm]}, thick, color=azulcos!81!gray] (axis cs:0,0) -- (axis cs:0.519,0.736) node[anchor=235, font=\scriptsize, inner sep=1pt] {\texttt{alcal\_ceniza}};
\draw[-{Stealth[length=2mm]}, thick, color=azulcos!12!gray] (axis cs:0,0) -- (axis cs:-0.308,0.157) node[anchor=333, font=\scriptsize, inner sep=1pt] {\texttt{magnesio}};
\draw[-{Stealth[length=2mm]}, thick, color=azulcos!76!gray] (axis cs:0,0) -- (axis cs:-0.856,0.176) node[anchor=348, font=\scriptsize, inner sep=1pt] {\texttt{fenoles\_tot}};
\draw[-{Stealth[length=2mm]}, thick, color=azulcos!87!gray] (axis cs:0,0) -- (axis cs:-0.917,0.181) node[anchor=349, font=\scriptsize, inner sep=1pt] {\texttt{flavonoides}};
\draw[-{Stealth[length=2mm]}, thick, color=azulcos!46!gray] (axis cs:0,0) -- (axis cs:0.648,0.205) node[anchor=198, font=\scriptsize, inner sep=1pt] {\texttt{fenoles\_noflav}};
\draw[-{Stealth[length=2mm]}, thick, color=azulcos!49!gray] (axis cs:0,0) -- (axis cs:-0.680,0.180) node[anchor=345, font=\scriptsize, inner sep=1pt] {\texttt{proantocianinas}};
\draw[-{Stealth[length=2mm]}, thick, color=azulcos!6!gray] (axis cs:0,0) -- (axis cs:0.192,-0.165) node[anchor=139, font=\scriptsize, inner sep=1pt] {\texttt{int\_color}};
\draw[-{Stealth[length=2mm]}, thick, color=azulcos!42!gray] (axis cs:0,0) -- (axis cs:-0.644,0.102) node[anchor=351, font=\scriptsize, inner sep=1pt] {\texttt{tono}};
\draw[-{Stealth[length=2mm]}, thick, color=azulcos!71!gray] (axis cs:0,0) -- (axis cs:-0.816,0.200) node[anchor=346, font=\scriptsize, inner sep=1pt] {\texttt{od280\_od315}};
\draw[-{Stealth[length=2mm]}, thick, color=azulcos!41!gray] (axis cs:0,0) -- (axis cs:-0.622,-0.152) node[anchor=14, font=\scriptsize, inner sep=1pt] {\texttt{prolina}};
\nextgroupplot[xlabel={CP2 (19.2\,\%)}, ylabel={CP3 (11.1\,\%)}, title={Plano (2,3)}]
\draw[black] (axis cs:0,0) circle[radius=1];
\draw[gray, dashed] (axis cs:-1.25,0) -- (axis cs:1.25,0); \draw[gray, dashed] (axis cs:0,-1.25) -- (axis cs:0,1.25);
\draw[-{Stealth[length=2mm]}, thick, color=azulcos!65!gray] (axis cs:0,0) -- (axis cs:-0.764,-0.249) node[anchor=18, font=\scriptsize, inner sep=1pt] {\texttt{alcohol}};
\draw[-{Stealth[length=2mm]}, thick, color=azulcos!14!gray] (axis cs:0,0) -- (axis cs:-0.355,0.107) node[anchor=343, font=\scriptsize, inner sep=1pt] {\texttt{ac\_malico}};
\draw[-{Stealth[length=2mm]}, thick, color=azulcos!82!gray] (axis cs:0,0) -- (axis cs:-0.499,0.753) node[anchor=304, font=\scriptsize, inner sep=1pt] {\texttt{ceniza}};
\draw[-{Stealth[length=2mm]}, thick, color=azulcos!54!gray] (axis cs:0,0) -- (axis cs:0.017,0.736) node[anchor=269, font=\scriptsize, inner sep=1pt] {\texttt{alcal\_ceniza}};
\draw[-{Stealth[length=2mm]}, thick, color=azulcos!25!gray] (axis cs:0,0) -- (axis cs:-0.473,0.157) node[anchor=342, font=\scriptsize, inner sep=1pt] {\texttt{magnesio}};
\draw[-{Stealth[length=2mm]}, thick, color=azulcos!4!gray] (axis cs:0,0) -- (axis cs:-0.103,0.176) node[anchor=300, font=\scriptsize, inner sep=1pt] {\texttt{fenoles\_tot}};
\draw[-{Stealth[length=2mm]}, thick, color=azulcos!3!gray] (axis cs:0,0) -- (axis cs:0.005,0.181) node[anchor=268, font=\scriptsize, inner sep=1pt] {\texttt{flavonoides}};
\draw[-{Stealth[length=2mm]}, thick, color=azulcos!4!gray] (axis cs:0,0) -- (axis cs:-0.045,0.205) node[anchor=283, font=\scriptsize, inner sep=1pt] {\texttt{fenoles\_noflav}};
\draw[-{Stealth[length=2mm]}, thick, color=azulcos!4!gray] (axis cs:0,0) -- (axis cs:-0.062,0.180) node[anchor=289, font=\scriptsize, inner sep=1pt] {\texttt{proantocianinas}};
\draw[-{Stealth[length=2mm]}, thick, color=azulcos!73!gray] (axis cs:0,0) -- (axis cs:-0.837,-0.165) node[anchor=11, font=\scriptsize, inner sep=1pt] {\texttt{int\_color}};
\draw[-{Stealth[length=2mm]}, thick, color=azulcos!21!gray] (axis cs:0,0) -- (axis cs:0.441,0.102) node[anchor=193, font=\scriptsize, inner sep=1pt] {\texttt{tono}};
\draw[-{Stealth[length=2mm]}, thick, color=azulcos!11!gray] (axis cs:0,0) -- (axis cs:0.260,0.200) node[anchor=218, font=\scriptsize, inner sep=1pt] {\texttt{od280\_od315}};
\draw[-{Stealth[length=2mm]}, thick, color=azulcos!36!gray] (axis cs:0,0) -- (axis cs:-0.577,-0.152) node[anchor=15, font=\scriptsize, inner sep=1pt] {\texttt{prolina}};
\end{groupplot}
\end{tikzpicture}
\caption{Círculos de correlaciones construidos con $R_{cv} = \operatorname{cor}(Z, F)$. La intensidad del color de cada flecha es proporcional a su $\cos^2$ en el plano (más oscuro = mejor representada).}
\label{fig:circ}
\end{figure}

\begin{figure}[H]
\centering
\begin{tikzpicture}
\begin{groupplot}[group style={group size=3 by 1, horizontal sep=2.6cm},
  width=4.6cm, height=7.2cm, ejes, xbar, /pgf/bar width=6pt, /pgf/bar shift=0pt, xmin=0, grid=none,
  y tick label style={font=\tiny}, xlabel={Contribución (\%)}]
\nextgroupplot[title={CP1 (36.2\,\%)}, symbolic y coords={ceniza,intcolor,magnesio,alcohol,alcalceniza,acmalico,prolina,tono,fenolesnoflav,proantocianinas,od280od315,fenolestot,flavonoides}, ytick={ceniza,intcolor,magnesio,alcohol,alcalceniza,acmalico,prolina,tono,fenolesnoflav,proantocianinas,od280od315,fenolestot,flavonoides}, yticklabels={\texttt{ceniza},\texttt{int\_color},\texttt{magnesio},\texttt{alcohol},\texttt{alcal\_ceniza},\texttt{ac\_malico},\texttt{prolina},\texttt{tono},\texttt{fenoles\_noflav},\texttt{proantocianinas},\texttt{od280\_od315},\texttt{fenoles\_tot},\texttt{flavonoides}}, xmax=22]
\addplot[fill=azul, draw=azul] coordinates {(8.22,prolina) (8.80,tono) (8.91,fenolesnoflav) (9.82,proantocianinas) (14.15,od280od315) (15.58,fenolestot) (17.89,flavonoides)};
\addplot[fill=gray!35, draw=gray!50] coordinates {(0.00,ceniza) (0.79,intcolor) (2.02,magnesio) (2.08,alcohol) (5.73,alcalceniza) (6.01,acmalico)};
\draw[rojo, dashed, thick] ({axis cs:7.69,ceniza}|-{rel axis cs:0,0}) -- ({axis cs:7.69,ceniza}|-{rel axis cs:0,1});
\nextgroupplot[title={CP2 (19.2\,\%)}, symbolic y coords={flavonoides,alcalceniza,fenolesnoflav,proantocianinas,fenolestot,od280od315,acmalico,tono,magnesio,ceniza,prolina,alcohol,intcolor}, ytick={flavonoides,alcalceniza,fenolesnoflav,proantocianinas,fenolestot,od280od315,acmalico,tono,magnesio,ceniza,prolina,alcohol,intcolor}, yticklabels={\texttt{flavonoides},\texttt{alcal\_ceniza},\texttt{fenoles\_noflav},\texttt{proantocianinas},\texttt{fenoles\_tot},\texttt{od280\_od315},\texttt{ac\_malico},\texttt{tono},\texttt{magnesio},\texttt{ceniza},\texttt{prolina},\texttt{alcohol},\texttt{int\_color}}, xmax=32]
\addplot[fill=azul, draw=azul] coordinates {(7.80,tono) (8.98,magnesio) (9.99,ceniza) (13.32,prolina) (23.39,alcohol) (28.09,intcolor)};
\addplot[fill=gray!35, draw=gray!50] coordinates {(0.00,flavonoides) (0.01,alcalceniza) (0.08,fenolesnoflav) (0.15,proantocianinas) (0.42,fenolestot) (2.71,od280od315) (5.06,acmalico)};
\draw[rojo, dashed, thick] ({axis cs:7.69,ceniza}|-{rel axis cs:0,0}) -- ({axis cs:7.69,ceniza}|-{rel axis cs:0,1});
\nextgroupplot[title={CP3 (11.1\,\%)}, symbolic y coords={tono,acmalico,prolina,magnesio,intcolor,fenolestot,proantocianinas,flavonoides,od280od315,fenolesnoflav,alcohol,alcalceniza,ceniza}, ytick={tono,acmalico,prolina,magnesio,intcolor,fenolestot,proantocianinas,flavonoides,od280od315,fenolesnoflav,alcohol,alcalceniza,ceniza}, yticklabels={\texttt{tono},\texttt{ac\_malico},\texttt{prolina},\texttt{magnesio},\texttt{int\_color},\texttt{fenoles\_tot},\texttt{proantocianinas},\texttt{flavonoides},\texttt{od280\_od315},\texttt{fenoles\_noflav},\texttt{alcohol},\texttt{alcal\_ceniza},\texttt{ceniza}}, xmax=44]
\addplot[fill=azul, draw=azul] coordinates {(37.46,alcalceniza) (39.22,ceniza)};
\addplot[fill=gray!35, draw=gray!50] coordinates {(0.73,tono) (0.79,acmalico) (1.61,prolina) (1.71,magnesio) (1.89,intcolor) (2.14,fenolestot) (2.23,proantocianinas) (2.27,flavonoides) (2.76,od280od315) (2.90,fenolesnoflav) (4.30,alcohol)};
\draw[rojo, dashed, thick] ({axis cs:7.69,ceniza}|-{rel axis cs:0,0}) -- ({axis cs:7.69,ceniza}|-{rel axis cs:0,1});
\end{groupplot}
\end{tikzpicture}
\caption{Contribución de las variables a cada componente retenida. La línea roja marca la contribución de referencia $1/p = 7.69\,\%$; en azul las variables que la superan.}
\label{fig:ctr}
\end{figure}

\begin{figure}[H]
\centering
\begin{tikzpicture}
\begin{groupplot}[group style={group size=3 by 1, horizontal sep=1.8cm},
  width=5.3cm, height=5.5cm, ejes, every axis plot/.append style={only marks}]
\nextgroupplot[xlabel={CP1 (36.2\,\%)}, ylabel={CP2 (19.2\,\%)}, title={Plano (1,2)}]
\addplot[mark=*, mark size=1.1pt, c1, opacity=.8, discard if not={cult}{1}] table[x=cpa, y=cpb]{\tablaF};
\addplot[mark=*, mark size=1.1pt, c2, opacity=.8, discard if not={cult}{2}] table[x=cpa, y=cpb]{\tablaF};
\addplot[mark=*, mark size=1.1pt, c3, opacity=.8, discard if not={cult}{3}] table[x=cpa, y=cpb]{\tablaF};
\addplot[c1, thick, mark=none, sharp plot, forget plot] coordinates {(-0.1156,-0.1328) (-0.1252,-0.1412) (-0.1351,-0.1491) (-0.1453,-0.1563) (-0.1557,-0.1630) (-0.1662,-0.1690) (-0.1768,-0.1742) (-0.1873,-0.1786) (-0.1977,-0.1823) (-0.2079,-0.1851) (-0.2178,-0.1871) (-0.2273,-0.1882) (-0.2364,-0.1884) (-0.2450,-0.1877) (-0.2531,-0.1861) (-0.2605,-0.1837) (-0.2673,-0.1805) (-0.2733,-0.1764) (-0.2786,-0.1715) (-0.2830,-0.1659) (-0.2866,-0.1595) (-0.2893,-0.1525) (-0.2911,-0.1449) (-0.2919,-0.1368) (-0.2919,-0.1281) (-0.2909,-0.1191) (-0.2891,-0.1096) (-0.2863,-0.0999) (-0.2827,-0.0900) (-0.2782,-0.0800) (-0.2729,-0.0699) (-0.2668,-0.0598) (-0.2600,-0.0498) (-0.2525,-0.0400) (-0.2444,-0.0304) (-0.2357,-0.0212) (-0.2266,-0.0123) (-0.2170,-0.0039) (-0.2071,0.0040) (-0.1969,0.0113) (-0.1865,0.0179) (-0.1760,0.0239) (-0.1654,0.0291) (-0.1549,0.0336) (-0.1445,0.0372) (-0.1343,0.0400) (-0.1244,0.0420) (-0.1149,0.0431) (-0.1058,0.0433) (-0.0972,0.0426) (-0.0891,0.0411) (-0.0817,0.0386) (-0.0749,0.0354) (-0.0689,0.0313) (-0.0636,0.0264) (-0.0592,0.0208) (-0.0556,0.0145) (-0.0529,0.0075) (-0.0511,-0.0002) (-0.0503,-0.0083) (-0.0503,-0.0170) (-0.0513,-0.0260) (-0.0531,-0.0355) (-0.0559,-0.0452) (-0.0595,-0.0551) (-0.0640,-0.0651) (-0.0693,-0.0752) (-0.0754,-0.0853) (-0.0822,-0.0953) (-0.0897,-0.1051) (-0.0978,-0.1147) (-0.1065,-0.1239) (-0.1156,-0.1328)};
\addplot[only marks, mark=diamond*, mark size=4pt, fill=c1, draw=black, forget plot] coordinates {(-0.1711,-0.0725)};
\addplot[c2, thick, mark=none, sharp plot, forget plot] coordinates {(0.0082,0.0002) (-0.0082,0.0000) (-0.0246,0.0007) (-0.0407,0.0024) (-0.0564,0.0049) (-0.0718,0.0084) (-0.0865,0.0127) (-0.1006,0.0179) (-0.1139,0.0239) (-0.1263,0.0306) (-0.1377,0.0380) (-0.1480,0.0461) (-0.1572,0.0548) (-0.1652,0.0640) (-0.1719,0.0736) (-0.1773,0.0836) (-0.1813,0.0940) (-0.1839,0.1045) (-0.1850,0.1152) (-0.1848,0.1259) (-0.1831,0.1367) (-0.1800,0.1473) (-0.1755,0.1577) (-0.1696,0.1679) (-0.1625,0.1777) (-0.1540,0.1871) (-0.1444,0.1961) (-0.1337,0.2045) (-0.1219,0.2122) (-0.1092,0.2193) (-0.0956,0.2257) (-0.0813,0.2312) (-0.0663,0.2360) (-0.0508,0.2399) (-0.0349,0.2429) (-0.0187,0.2450) (-0.0023,0.2461) (0.0141,0.2464) (0.0304,0.2457) (0.0465,0.2440) (0.0623,0.2415) (0.0776,0.2380) (0.0924,0.2337) (0.1064,0.2285) (0.1197,0.2225) (0.1321,0.2158) (0.1435,0.2083) (0.1539,0.2003) (0.1631,0.1916) (0.1711,0.1824) (0.1778,0.1728) (0.1831,0.1627) (0.1871,0.1524) (0.1897,0.1419) (0.1909,0.1312) (0.1906,0.1204) (0.1889,0.1097) (0.1858,0.0991) (0.1813,0.0887) (0.1755,0.0785) (0.1683,0.0686) (0.1599,0.0592) (0.1503,0.0503) (0.1395,0.0419) (0.1278,0.0341) (0.1150,0.0271) (0.1014,0.0207) (0.0871,0.0151) (0.0721,0.0104) (0.0566,0.0065) (0.0407,0.0035) (0.0245,0.0014) (0.0082,0.0002)};
\addplot[only marks, mark=diamond*, mark size=4pt, fill=c2, draw=black, forget plot] coordinates {(0.0029,0.1232)};
\addplot[c3, thick, mark=none, sharp plot, forget plot] coordinates {(0.1186,-0.0664) (0.1226,-0.0545) (0.1273,-0.0428) (0.1325,-0.0315) (0.1383,-0.0207) (0.1446,-0.0104) (0.1514,-0.0008) (0.1586,0.0081) (0.1662,0.0163) (0.1740,0.0236) (0.1821,0.0301) (0.1904,0.0356) (0.1988,0.0401) (0.2073,0.0436) (0.2157,0.0461) (0.2241,0.0475) (0.2323,0.0479) (0.2404,0.0471) (0.2481,0.0453) (0.2556,0.0425) (0.2627,0.0386) (0.2693,0.0337) (0.2755,0.0279) (0.2811,0.0211) (0.2862,0.0135) (0.2906,0.0050) (0.2944,-0.0042) (0.2976,-0.0140) (0.3000,-0.0245) (0.3017,-0.0355) (0.3027,-0.0469) (0.3030,-0.0587) (0.3025,-0.0707) (0.3013,-0.0829) (0.2994,-0.0952) (0.2967,-0.1075) (0.2934,-0.1197) (0.2893,-0.1316) (0.2847,-0.1433) (0.2795,-0.1546) (0.2737,-0.1654) (0.2674,-0.1756) (0.2606,-0.1853) (0.2534,-0.1942) (0.2458,-0.2024) (0.2380,-0.2097) (0.2298,-0.2162) (0.2216,-0.2217) (0.2132,-0.2262) (0.2047,-0.2297) (0.1962,-0.2322) (0.1879,-0.2336) (0.1796,-0.2339) (0.1716,-0.2332) (0.1638,-0.2314) (0.1564,-0.2286) (0.1493,-0.2247) (0.1426,-0.2198) (0.1365,-0.2140) (0.1308,-0.2072) (0.1258,-0.1996) (0.1213,-0.1911) (0.1175,-0.1819) (0.1144,-0.1721) (0.1120,-0.1616) (0.1102,-0.1506) (0.1092,-0.1392) (0.1090,-0.1274) (0.1095,-0.1153) (0.1107,-0.1031) (0.1126,-0.0908) (0.1153,-0.0786) (0.1186,-0.0664)};
\addplot[only marks, mark=diamond*, mark size=4pt, fill=c3, draw=black, forget plot] coordinates {(0.2060,-0.0930)};
\nextgroupplot[xlabel={CP1 (36.2\,\%)}, ylabel={CP3 (11.1\,\%)}, title={Plano (1,3)}]
\addplot[mark=*, mark size=1.1pt, c1, opacity=.8, discard if not={cult}{1}] table[x=cpa, y=cpc]{\tablaF};
\addplot[mark=*, mark size=1.1pt, c2, opacity=.8, discard if not={cult}{2}] table[x=cpa, y=cpc]{\tablaF};
\addplot[mark=*, mark size=1.1pt, c3, opacity=.8, discard if not={cult}{3}] table[x=cpa, y=cpc]{\tablaF};
\addplot[c1, thick, mark=none, sharp plot, forget plot] coordinates {(-0.2610,0.0401) (-0.2536,0.0521) (-0.2455,0.0636) (-0.2370,0.0745) (-0.2279,0.0847) (-0.2183,0.0943) (-0.2085,0.1030) (-0.1983,0.1108) (-0.1879,0.1177) (-0.1774,0.1236) (-0.1669,0.1285) (-0.1563,0.1323) (-0.1459,0.1350) (-0.1357,0.1366) (-0.1258,0.1370) (-0.1162,0.1363) (-0.1070,0.1345) (-0.0983,0.1316) (-0.0902,0.1276) (-0.0826,0.1225) (-0.0758,0.1165) (-0.0697,0.1094) (-0.0643,0.1014) (-0.0598,0.0925) (-0.0561,0.0829) (-0.0533,0.0725) (-0.0513,0.0615) (-0.0503,0.0499) (-0.0502,0.0378) (-0.0511,0.0254) (-0.0528,0.0127) (-0.0555,-0.0002) (-0.0590,-0.0132) (-0.0634,-0.0262) (-0.0685,-0.0390) (-0.0745,-0.0517) (-0.0812,-0.0641) (-0.0886,-0.0760) (-0.0967,-0.0875) (-0.1052,-0.0984) (-0.1143,-0.1087) (-0.1239,-0.1182) (-0.1337,-0.1269) (-0.1439,-0.1347) (-0.1543,-0.1416) (-0.1648,-0.1475) (-0.1753,-0.1524) (-0.1858,-0.1562) (-0.1963,-0.1589) (-0.2065,-0.1605) (-0.2164,-0.1609) (-0.2260,-0.1603) (-0.2352,-0.1585) (-0.2439,-0.1555) (-0.2520,-0.1515) (-0.2596,-0.1465) (-0.2664,-0.1404) (-0.2725,-0.1333) (-0.2779,-0.1253) (-0.2824,-0.1164) (-0.2861,-0.1068) (-0.2889,-0.0964) (-0.2909,-0.0854) (-0.2919,-0.0738) (-0.2920,-0.0618) (-0.2911,-0.0493) (-0.2894,-0.0366) (-0.2867,-0.0237) (-0.2832,-0.0107) (-0.2788,0.0022) (-0.2736,0.0151) (-0.2677,0.0278) (-0.2610,0.0401)};
\addplot[only marks, mark=diamond*, mark size=4pt, fill=c1, draw=black, forget plot] coordinates {(-0.1711,-0.0120)};
\addplot[c2, thick, mark=none, sharp plot, forget plot] coordinates {(-0.1777,-0.0237) (-0.1816,-0.0044) (-0.1840,0.0149) (-0.1851,0.0344) (-0.1847,0.0537) (-0.1829,0.0727) (-0.1797,0.0914) (-0.1750,0.1095) (-0.1691,0.1269) (-0.1618,0.1435) (-0.1533,0.1591) (-0.1435,0.1737) (-0.1327,0.1872) (-0.1208,0.1993) (-0.1080,0.2101) (-0.0944,0.2194) (-0.0800,0.2272) (-0.0649,0.2335) (-0.0494,0.2381) (-0.0335,0.2410) (-0.0172,0.2423) (-0.0009,0.2418) (0.0155,0.2397) (0.0318,0.2359) (0.0479,0.2304) (0.0636,0.2234) (0.0789,0.2148) (0.0936,0.2047) (0.1076,0.1932) (0.1209,0.1803) (0.1332,0.1663) (0.1445,0.1511) (0.1547,0.1350) (0.1638,0.1179) (0.1717,0.1001) (0.1783,0.0817) (0.1835,0.0629) (0.1874,0.0437) (0.1899,0.0243) (0.1909,0.0049) (0.1905,-0.0145) (0.1887,-0.0335) (0.1855,-0.0522) (0.1809,-0.0703) (0.1749,-0.0877) (0.1676,-0.1043) (0.1591,-0.1199) (0.1494,-0.1345) (0.1385,-0.1480) (0.1267,-0.1601) (0.1139,-0.1709) (0.1002,-0.1802) (0.0858,-0.1880) (0.0708,-0.1943) (0.0552,-0.1989) (0.0393,-0.2018) (0.0231,-0.2030) (0.0067,-0.2026) (-0.0097,-0.2005) (-0.0260,-0.1967) (-0.0421,-0.1912) (-0.0578,-0.1842) (-0.0731,-0.1755) (-0.0878,-0.1655) (-0.1018,-0.1540) (-0.1150,-0.1411) (-0.1273,-0.1271) (-0.1386,-0.1119) (-0.1489,-0.0958) (-0.1580,-0.0787) (-0.1659,-0.0609) (-0.1724,-0.0425) (-0.1777,-0.0237)};
\addplot[only marks, mark=diamond*, mark size=4pt, fill=c2, draw=black, forget plot] coordinates {(0.0029,0.0196)};
\addplot[c3, thick, mark=none, sharp plot, forget plot] coordinates {(0.1112,-0.0290) (0.1098,-0.0174) (0.1091,-0.0057) (0.1091,0.0059) (0.1098,0.0174) (0.1113,0.0286) (0.1136,0.0395) (0.1165,0.0500) (0.1201,0.0600) (0.1243,0.0694) (0.1292,0.0782) (0.1347,0.0863) (0.1407,0.0936) (0.1472,0.1001) (0.1541,0.1057) (0.1615,0.1105) (0.1691,0.1142) (0.1771,0.1170) (0.1853,0.1188) (0.1936,0.1196) (0.2021,0.1193) (0.2105,0.1181) (0.2190,0.1158) (0.2273,0.1126) (0.2355,0.1083) (0.2434,0.1032) (0.2511,0.0971) (0.2584,0.0902) (0.2653,0.0825) (0.2718,0.0741) (0.2777,0.0650) (0.2831,0.0553) (0.2880,0.0451) (0.2922,0.0344) (0.2957,0.0233) (0.2986,0.0120) (0.3008,0.0004) (0.3022,-0.0112) (0.3029,-0.0229) (0.3029,-0.0345) (0.3021,-0.0460) (0.3006,-0.0572) (0.2984,-0.0681) (0.2955,-0.0786) (0.2919,-0.0886) (0.2876,-0.0980) (0.2828,-0.1068) (0.2773,-0.1149) (0.2713,-0.1222) (0.2648,-0.1287) (0.2579,-0.1343) (0.2505,-0.1391) (0.2428,-0.1428) (0.2349,-0.1456) (0.2267,-0.1474) (0.2183,-0.1482) (0.2099,-0.1479) (0.2014,-0.1467) (0.1930,-0.1444) (0.1847,-0.1412) (0.1765,-0.1369) (0.1686,-0.1318) (0.1609,-0.1257) (0.1536,-0.1188) (0.1467,-0.1112) (0.1402,-0.1027) (0.1343,-0.0936) (0.1288,-0.0839) (0.1240,-0.0737) (0.1198,-0.0630) (0.1162,-0.0519) (0.1134,-0.0406) (0.1112,-0.0290)};
\addplot[only marks, mark=diamond*, mark size=4pt, fill=c3, draw=black, forget plot] coordinates {(0.2060,-0.0143)};
\nextgroupplot[xlabel={CP2 (19.2\,\%)}, ylabel={CP3 (11.1\,\%)}, title={Plano (2,3)}]
\addplot[mark=*, mark size=1.1pt, c1, opacity=.8, discard if not={cult}{1}] table[x=cpb, y=cpc]{\tablaF};
\addplot[mark=*, mark size=1.1pt, c2, opacity=.8, discard if not={cult}{2}] table[x=cpb, y=cpc]{\tablaF};
\addplot[mark=*, mark size=1.1pt, c3, opacity=.8, discard if not={cult}{3}] table[x=cpb, y=cpc]{\tablaF};
\addplot[c1, thick, mark=none, sharp plot, forget plot] coordinates {(-0.1810,-0.0422) (-0.1842,-0.0293) (-0.1864,-0.0164) (-0.1879,-0.0034) (-0.1884,0.0095) (-0.1880,0.0223) (-0.1868,0.0348) (-0.1847,0.0470) (-0.1818,0.0587) (-0.1780,0.0698) (-0.1734,0.0804) (-0.1681,0.0902) (-0.1620,0.0993) (-0.1552,0.1075) (-0.1478,0.1148) (-0.1399,0.1211) (-0.1314,0.1265) (-0.1225,0.1307) (-0.1132,0.1339) (-0.1036,0.1360) (-0.0937,0.1370) (-0.0837,0.1368) (-0.0736,0.1355) (-0.0635,0.1330) (-0.0535,0.1295) (-0.0436,0.1249) (-0.0339,0.1192) (-0.0246,0.1126) (-0.0155,0.1050) (-0.0070,0.0965) (0.0011,0.0872) (0.0086,0.0771) (0.0155,0.0663) (0.0217,0.0550) (0.0272,0.0431) (0.0320,0.0309) (0.0359,0.0183) (0.0391,0.0054) (0.0414,-0.0075) (0.0428,-0.0205) (0.0433,-0.0335) (0.0430,-0.0462) (0.0417,-0.0587) (0.0396,-0.0709) (0.0367,-0.0826) (0.0329,-0.0938) (0.0283,-0.1043) (0.0230,-0.1141) (0.0169,-0.1232) (0.0101,-0.1314) (0.0027,-0.1387) (-0.0052,-0.1451) (-0.0137,-0.1504) (-0.0226,-0.1547) (-0.0319,-0.1578) (-0.0415,-0.1599) (-0.0514,-0.1609) (-0.0614,-0.1607) (-0.0715,-0.1594) (-0.0816,-0.1569) (-0.0916,-0.1534) (-0.1015,-0.1488) (-0.1112,-0.1431) (-0.1205,-0.1365) (-0.1295,-0.1289) (-0.1381,-0.1204) (-0.1462,-0.1111) (-0.1537,-0.1010) (-0.1606,-0.0902) (-0.1668,-0.0789) (-0.1723,-0.0670) (-0.1771,-0.0548) (-0.1810,-0.0422)};
\addplot[only marks, mark=diamond*, mark size=4pt, fill=c1, draw=black, forget plot] coordinates {(-0.0725,-0.0120)};
\addplot[c2, thick, mark=none, sharp plot, forget plot] coordinates {(0.0169,-0.0104) (0.0119,0.0089) (0.0077,0.0284) (0.0044,0.0477) (0.0020,0.0669) (0.0005,0.0856) (-0.0000,0.1039) (0.0004,0.1216) (0.0017,0.1384) (0.0039,0.1544) (0.0071,0.1693) (0.0111,0.1831) (0.0160,0.1957) (0.0218,0.2069) (0.0282,0.2167) (0.0354,0.2250) (0.0433,0.2317) (0.0518,0.2368) (0.0608,0.2403) (0.0703,0.2421) (0.0802,0.2421) (0.0904,0.2405) (0.1009,0.2372) (0.1116,0.2323) (0.1223,0.2257) (0.1330,0.2176) (0.1437,0.2080) (0.1542,0.1969) (0.1645,0.1845) (0.1745,0.1708) (0.1840,0.1559) (0.1931,0.1401) (0.2017,0.1233) (0.2097,0.1057) (0.2170,0.0875) (0.2236,0.0688) (0.2294,0.0496) (0.2345,0.0303) (0.2387,0.0109) (0.2420,-0.0085) (0.2444,-0.0276) (0.2458,-0.0464) (0.2464,-0.0647) (0.2460,-0.0824) (0.2447,-0.0992) (0.2424,-0.1152) (0.2393,-0.1301) (0.2352,-0.1439) (0.2303,-0.1565) (0.2246,-0.1677) (0.2181,-0.1775) (0.2109,-0.1858) (0.2031,-0.1925) (0.1946,-0.1976) (0.1855,-0.2011) (0.1761,-0.2028) (0.1661,-0.2029) (0.1559,-0.2013) (0.1454,-0.1980) (0.1348,-0.1931) (0.1241,-0.1865) (0.1133,-0.1784) (0.1027,-0.1687) (0.0921,-0.1577) (0.0819,-0.1452) (0.0719,-0.1316) (0.0623,-0.1167) (0.0532,-0.1009) (0.0447,-0.0841) (0.0367,-0.0665) (0.0294,-0.0483) (0.0228,-0.0295) (0.0169,-0.0104)};
\addplot[only marks, mark=diamond*, mark size=4pt, fill=c2, draw=black, forget plot] coordinates {(0.1232,0.0196)};
\addplot[c3, thick, mark=none, sharp plot, forget plot] coordinates {(-0.1674,-0.1032) (-0.1776,-0.0941) (-0.1871,-0.0844) (-0.1959,-0.0742) (-0.2039,-0.0636) (-0.2110,-0.0525) (-0.2173,-0.0412) (-0.2226,-0.0296) (-0.2270,-0.0180) (-0.2303,-0.0063) (-0.2325,0.0053) (-0.2337,0.0168) (-0.2339,0.0280) (-0.2329,0.0389) (-0.2309,0.0494) (-0.2279,0.0595) (-0.2238,0.0689) (-0.2187,0.0777) (-0.2127,0.0859) (-0.2058,0.0932) (-0.1980,0.0998) (-0.1894,0.1055) (-0.1800,0.1102) (-0.1700,0.1141) (-0.1595,0.1169) (-0.1484,0.1187) (-0.1369,0.1196) (-0.1250,0.1194) (-0.1130,0.1182) (-0.1007,0.1160) (-0.0884,0.1128) (-0.0762,0.1086) (-0.0640,0.1035) (-0.0521,0.0975) (-0.0405,0.0906) (-0.0293,0.0830) (-0.0186,0.0746) (-0.0085,0.0655) (0.0010,0.0558) (0.0098,0.0456) (0.0178,0.0350) (0.0250,0.0239) (0.0312,0.0126) (0.0365,0.0010) (0.0409,-0.0106) (0.0442,-0.0223) (0.0465,-0.0339) (0.0477,-0.0454) (0.0478,-0.0566) (0.0469,-0.0675) (0.0449,-0.0780) (0.0418,-0.0881) (0.0377,-0.0975) (0.0326,-0.1064) (0.0266,-0.1145) (0.0197,-0.1218) (0.0119,-0.1284) (0.0033,-0.1341) (-0.0060,-0.1388) (-0.0160,-0.1427) (-0.0266,-0.1455) (-0.0377,-0.1473) (-0.0492,-0.1482) (-0.0610,-0.1480) (-0.0731,-0.1468) (-0.0854,-0.1446) (-0.0976,-0.1414) (-0.1099,-0.1372) (-0.1220,-0.1321) (-0.1339,-0.1261) (-0.1455,-0.1192) (-0.1567,-0.1116) (-0.1674,-0.1032)};
\addplot[only marks, mark=diamond*, mark size=4pt, fill=c3, draw=black, forget plot] coordinates {(-0.0930,-0.0143)};
\end{groupplot}
\end{tikzpicture}

\ref*{leyCult}
\caption{Vinos por cultivar en los planos (1,2), (1,3) y (2,3), con elipses de 2 desviaciones estándar y centroides ($\blacklozenge$).}
\label{fig:planos}
\end{figure}

\begin{figure}[H]
\centering
\begin{tikzpicture}
\begin{axis}[ejes, width=12cm, height=10cm, axis equal,
  xlabel={CP1 (36.2\,\%)}, ylabel={CP2 (19.2\,\%)},
  legend style={font=\footnotesize, at={(0.02,0.98)}, anchor=north west},
  title={Biplot -- vinos y variables químicas (plano 1-2)}]
\addplot[only marks, mark=*, mark size=1.3pt, c1, opacity=.6, discard if not={cult}{1}] table[x=cpa, y=cpb]{\tablaF};
\addplot[only marks, mark=*, mark size=1.3pt, c2, opacity=.6, discard if not={cult}{2}] table[x=cpa, y=cpb]{\tablaF};
\addplot[only marks, mark=*, mark size=1.3pt, c3, opacity=.6, discard if not={cult}{3}] table[x=cpa, y=cpb]{\tablaF};
\legend{Cultivar 1, Cultivar 2, Cultivar 3}
\draw[gray, dashed] (axis cs:-0.4,0) -- (axis cs:0.4,0);
\draw[gray, dashed] (axis cs:0,-0.35) -- (axis cs:0,0.35);
\draw[-{Stealth[length=2mm]}, thick, gray!30!black] (axis cs:0,0) -- (axis cs:-0.0993,-0.2423) node[anchor=68, font=\scriptsize, inner sep=1pt, fill=white, fill opacity=.7, text opacity=1] {\texttt{alcohol}};
\draw[-{Stealth[length=2mm]}, thick, gray!30!black] (axis cs:0,0) -- (axis cs:0.1687,-0.1127) node[anchor=146, font=\scriptsize, inner sep=1pt, fill=white, fill opacity=.7, text opacity=1] {\texttt{ac\_malico}};
\draw[-{Stealth[length=2mm]}, thick, gray!30!black] (axis cs:0,0) -- (axis cs:0.0014,-0.1584) node[anchor=91, font=\scriptsize, inner sep=1pt, fill=white, fill opacity=.7, text opacity=1] {\texttt{ceniza}};
\draw[-{Stealth[length=2mm]}, thick, gray!30!black] (axis cs:0,0) -- (axis cs:0.1646,0.0053) node[anchor=182, font=\scriptsize, inner sep=1pt, fill=white, fill opacity=.7, text opacity=1] {\texttt{alcal\_ceniza}};
\draw[-{Stealth[length=2mm]}, thick, gray!30!black] (axis cs:0,0) -- (axis cs:-0.0977,-0.1501) node[anchor=57, font=\scriptsize, inner sep=1pt, fill=white, fill opacity=.7, text opacity=1] {\texttt{magnesio}};
\draw[-{Stealth[length=2mm]}, thick, gray!30!black] (axis cs:0,0) -- (axis cs:-0.2715,-0.0326) node[anchor=7, font=\scriptsize, inner sep=1pt, fill=white, fill opacity=.7, text opacity=1] {\texttt{fenoles\_tot}};
\draw[-{Stealth[length=2mm]}, thick, gray!30!black] (axis cs:0,0) -- (axis cs:-0.2909,0.0017) node[anchor=360, font=\scriptsize, inner sep=1pt, fill=white, fill opacity=.7, text opacity=1] {\texttt{flavonoides}};
\draw[-{Stealth[length=2mm]}, thick, gray!30!black] (axis cs:0,0) -- (axis cs:0.2054,-0.0144) node[anchor=176, font=\scriptsize, inner sep=1pt, fill=white, fill opacity=.7, text opacity=1] {\texttt{fenoles\_noflav}};
\draw[-{Stealth[length=2mm]}, thick, gray!30!black] (axis cs:0,0) -- (axis cs:-0.2156,-0.0197) node[anchor=5, font=\scriptsize, inner sep=1pt, fill=white, fill opacity=.7, text opacity=1] {\texttt{proantocianinas}};
\draw[-{Stealth[length=2mm]}, thick, gray!30!black] (axis cs:0,0) -- (axis cs:0.0610,-0.2656) node[anchor=103, font=\scriptsize, inner sep=1pt, fill=white, fill opacity=.7, text opacity=1] {\texttt{int\_color}};
\draw[-{Stealth[length=2mm]}, thick, gray!30!black] (axis cs:0,0) -- (axis cs:-0.2041,0.1399) node[anchor=326, font=\scriptsize, inner sep=1pt, fill=white, fill opacity=.7, text opacity=1] {\texttt{tono}};
\draw[-{Stealth[length=2mm]}, thick, gray!30!black] (axis cs:0,0) -- (axis cs:-0.2588,0.0824) node[anchor=342, font=\scriptsize, inner sep=1pt, fill=white, fill opacity=.7, text opacity=1] {\texttt{od280\_od315}};
\draw[-{Stealth[length=2mm]}, thick, gray!30!black] (axis cs:0,0) -- (axis cs:-0.1973,-0.1828) node[anchor=43, font=\scriptsize, inner sep=1pt, fill=white, fill opacity=.7, text opacity=1] {\texttt{prolina}};
\end{axis}
\end{tikzpicture}
\caption{Biplot del primer plano: individuos y variables (flechas reescaladas a la escala de las coordenadas).}
\label{fig:biplot}
\end{figure}

\paragraph{Lectura de los gráficos.} En el primer plano ($\approx 55.4\,\%$ de la inercia) los tres cultivares forman tres nubes casi separadas en forma de arco. CP1 separa el cultivar 1 (izquierda) del cultivar 3 (derecha), y CP2 aísla al cultivar 2 (arriba). En los planos (1,3) y (2,3) la CP3 no separa los cultivares: las elipses se superponen verticalmente.

\paragraph{Reflexión.} El biplot muestra simultáneamente observaciones y variables, lo que facilita ver qué características químicas explican la posición de cada grupo de vinos, algo que por separado el gráfico de individuos y el círculo de correlaciones no muestran. Su limitación es que, al visualizar solo dos dimensiones, se pierde la información de las componentes de orden superior.

% =====================================================================
\section{Interpretación sustantiva de las componentes}

\begin{table}[H]
\centering\small
\caption{Cargas ($Q$), correlaciones ($R_{cv}$), $\cos^2$ y contribuciones de las variables a cada componente retenida.}
\label{tab:load}
\begin{tabular}{lrrrrcc}
\toprule
Variable & Carga ($Q$) & $R_{cv}$ & $\cos^2$ & ctr (\%) & ctr $> 1/p$ & $\cos^2 > 0.3$ \\
\midrule
\multicolumn{7}{l}{\textbf{CP1 (36.2\,\%)}} \\
\texttt{flavonoides} & -0.423 & -0.917 & 0.842 & 17.887 & $\checkmark$ & $\checkmark$ \\
\texttt{fenoles\_tot} & -0.395 & -0.856 & 0.733 & 15.576 & $\checkmark$ & $\checkmark$ \\
\texttt{od280\_od315} & -0.376 & -0.816 & 0.666 & 14.150 & $\checkmark$ & $\checkmark$ \\
\texttt{proantocianinas} & -0.313 & -0.680 & 0.462 & 9.824 & $\checkmark$ & $\checkmark$ \\
\texttt{fenoles\_noflav} & 0.299 & 0.648 & 0.419 & 8.912 & $\checkmark$ & $\checkmark$ \\
\texttt{tono} & -0.297 & -0.644 & 0.414 & 8.804 & $\checkmark$ & $\checkmark$ \\
\texttt{prolina} & -0.287 & -0.622 & 0.387 & 8.223 & $\checkmark$ & $\checkmark$ \\
\texttt{ac\_malico} & 0.245 & 0.532 & 0.283 & 6.012 &  &  \\
\texttt{alcal\_ceniza} & 0.239 & 0.519 & 0.270 & 5.727 &  &  \\
\texttt{alcohol} & -0.144 & -0.313 & 0.098 & 2.083 &  &  \\
\texttt{magnesio} & -0.142 & -0.308 & 0.095 & 2.016 &  &  \\
\texttt{int\_color} & 0.089 & 0.192 & 0.037 & 0.785 &  &  \\
\texttt{ceniza} & 0.002 & 0.004 & 0.000 & 0.000 &  &  \\
\midrule
\multicolumn{7}{l}{\textbf{CP2 (19.2\,\%)}} \\
\texttt{int\_color} & -0.530 & -0.837 & 0.701 & 28.090 & $\checkmark$ & $\checkmark$ \\
\texttt{alcohol} & -0.484 & -0.764 & 0.584 & 23.392 & $\checkmark$ & $\checkmark$ \\
\texttt{prolina} & -0.365 & -0.577 & 0.332 & 13.315 & $\checkmark$ & $\checkmark$ \\
\texttt{ceniza} & -0.316 & -0.499 & 0.249 & 9.990 & $\checkmark$ &  \\
\texttt{magnesio} & -0.300 & -0.473 & 0.224 & 8.978 & $\checkmark$ &  \\
\texttt{tono} & 0.279 & 0.441 & 0.195 & 7.797 & $\checkmark$ &  \\
\texttt{ac\_malico} & -0.225 & -0.355 & 0.126 & 5.059 &  &  \\
\texttt{od280\_od315} & 0.164 & 0.260 & 0.068 & 2.706 &  &  \\
\texttt{fenoles\_tot} & -0.065 & -0.103 & 0.011 & 0.423 &  &  \\
\texttt{proantocianinas} & -0.039 & -0.062 & 0.004 & 0.154 &  &  \\
\texttt{fenoles\_noflav} & -0.029 & -0.045 & 0.002 & 0.083 &  &  \\
\texttt{alcal\_ceniza} & 0.011 & 0.017 & 0.000 & 0.011 &  &  \\
\texttt{flavonoides} & 0.003 & 0.005 & 0.000 & 0.001 &  &  \\
\midrule
\multicolumn{7}{l}{\textbf{CP3 (11.1\,\%)}} \\
\texttt{ceniza} & 0.626 & 0.753 & 0.567 & 39.216 & $\checkmark$ & $\checkmark$ \\
\texttt{alcal\_ceniza} & 0.612 & 0.736 & 0.542 & 37.464 & $\checkmark$ & $\checkmark$ \\
\texttt{alcohol} & -0.207 & -0.249 & 0.062 & 4.301 &  &  \\
\texttt{fenoles\_noflav} & 0.170 & 0.205 & 0.042 & 2.903 &  &  \\
\texttt{od280\_od315} & 0.166 & 0.200 & 0.040 & 2.756 &  &  \\
\texttt{flavonoides} & 0.151 & 0.181 & 0.033 & 2.271 &  &  \\
\texttt{proantocianinas} & 0.149 & 0.180 & 0.032 & 2.234 &  &  \\
\texttt{fenoles\_tot} & 0.146 & 0.176 & 0.031 & 2.137 &  &  \\
\texttt{int\_color} & -0.137 & -0.165 & 0.027 & 1.885 &  &  \\
\texttt{magnesio} & 0.131 & 0.157 & 0.025 & 1.710 &  &  \\
\texttt{prolina} & -0.127 & -0.152 & 0.023 & 1.606 &  &  \\
\texttt{ac\_malico} & 0.089 & 0.107 & 0.011 & 0.792 &  &  \\
\texttt{tono} & 0.085 & 0.102 & 0.011 & 0.726 &  &  \\
\bottomrule
\end{tabular}
\end{table}

\begin{table}[H]
\centering\small
\caption{Lectura técnica de cada eje (magnitudes grandes de $R_{cv}$ agrupadas por signo; el signo es arbitrario).}
\begin{tabularx}{\textwidth}{>{\bfseries}l >{\raggedright\arraybackslash}X >{\raggedright\arraybackslash}X >{\raggedright\arraybackslash}p{3.2cm}}
\toprule
Eje & Lado negativo & Lado positivo & Cultivar en cada lado \\
\midrule
CP1 & \var{flavonoides} ($-0.92$), \var{fenoles\_tot} ($-0.86$), \var{od280\_od315} ($-0.82$), \var{proantocianinas} ($-0.68$), \var{tono} ($-0.64$), \var{prolina} ($-0.62$)
    & \var{fenoles\_noflav} ($0.65$), \var{ac\_malico} ($0.53$), \var{alcal\_ceniza} ($0.52$)
    & Negativo: cultivar 1. Positivo: cultivar 3 \\
CP2 & \var{int\_color} ($-0.84$), \var{alcohol} ($-0.76$), \var{prolina} ($-0.58$), \var{ceniza} ($-0.50$), \var{magnesio} ($-0.47$)
    & \var{tono} ($0.44$)
    & Negativo: cultivares 1 y 3. Positivo: cultivar 2 \\
CP3 & -- & \var{ceniza} ($0.75$), \var{alcal\_ceniza} ($0.74$) & No separa cultivares \\
\bottomrule
\end{tabularx}
\end{table}

\paragraph{CP1 (36.2\,\%): eje de riqueza fenólica y calidad antioxidante.} Es la dimensión más importante: concentra más de un tercio de la variabilidad total. Mide la intensidad del perfil antioxidante y la complejidad química derivada de los compuestos fenólicos. Con el signo que entrega la DVS, los valores \emph{negativos} corresponden a vinos con alta concentración de flavonoides, fenoles totales y proantocianidinas (cultivar 1), típicamente asociados a vinos de mayor cuerpo y estructura; los positivos, a vinos con más fenoles no flavonoides y ácido málico (cultivar 3).

\paragraph{CP2 (19.2\,\%): eje de perfil cromático, alcohol y madurez.} Capta casi una quinta parte de la información. Lo dominan la intensidad de color, el alcohol y la prolina, opuestos al tono. Un vino en los extremos de este eje se diferencia principalmente por la intensidad u oscuridad de su color y su grado alcohólico, independientemente de su cantidad de fenoles totales. Los vinos del cultivar 2 (más claros y menos alcohólicos) quedan en el lado positivo.

\paragraph{CP3 (11.1\,\%): eje de composición mineral.} Aporta una décima parte de la variabilidad y recoge casi exclusivamente \var{ceniza} y \var{alcal\_ceniza}, variables inorgánicas que no fueron explicadas por los dos primeros ejes. Describe el fondo mineral del vino y no discrimina entre cultivares.

% =====================================================================
\section{Proyección de las observaciones, patrones y atípicos}

\begin{table}[H]
\centering\small
\caption{Vinos mejor representados ($\cos^2$) y que más contribuyen (ctr) a cada componente. Contribución promedio de un vino: $100/n = 0.562\,\%$.}
\begin{tabular}{llrclrc}
\toprule
 & \multicolumn{3}{c}{Mejor representados} & \multicolumn{3}{c}{Mayor contribución} \\
\cmidrule(lr){2-4}\cmidrule(lr){5-7}
CP & Vino & $\cos^2$ & Cult. & Vino & ctr (\%) & Cult. \\
\midrule
\multirow{5}{*}{CP1} & Vino\_171 & 0.820 & 3 & Vino\_015 & 2.221 & 1 \\
 & Vino\_137 & 0.780 & 3 & Vino\_147 & 2.188 & 3 \\
 & Vino\_166 & 0.772 & 3 & Vino\_138 & 1.850 & 3 \\
 & Vino\_148 & 0.765 & 3 & Vino\_137 & 1.831 & 3 \\
 & Vino\_036 & 0.751 & 1 & Vino\_004 & 1.685 & 1 \\
\midrule
\multirow{5}{*}{CP2} & Vino\_081 & 0.830 & 2 & Vino\_116 & 3.373 & 2 \\
 & Vino\_117 & 0.805 & 2 & Vino\_159 & 2.780 & 3 \\
 & Vino\_104 & 0.737 & 2 & Vino\_081 & 2.563 & 2 \\
 & Vino\_107 & 0.699 & 2 & Vino\_060 & 2.125 & 2 \\
 & Vino\_098 & 0.670 & 2 & Vino\_117 & 1.791 & 2 \\
\midrule
\multirow{5}{*}{CP3} & Vino\_026 & 0.769 & 1 & Vino\_122 & 11.101 & 2 \\
 & Vino\_122 & 0.751 & 2 & Vino\_060 & 8.167 & 2 \\
 & Vino\_128 & 0.660 & 2 & Vino\_026 & 5.672 & 1 \\
 & Vino\_123 & 0.633 & 2 & Vino\_074 & 4.424 & 2 \\
 & Vino\_060 & 0.613 & 2 & Vino\_128 & 4.389 & 2 \\

\bottomrule
\end{tabular}
\end{table}

\begin{table}[H]
\centering\small
\caption{Coordenadas de tres vinos concretos y su calidad de representación en el espacio de 3 componentes.}
\begin{tabular}{lcrrrr}
\toprule
Vino & Cultivar & CP1 & CP2 & CP3 & $\cos^2$ (3 CP) \\
\midrule
Vino\_159 & 3 & 0.079 & -0.263 & 0.087 & 0.498 \\
Vino\_074 & 2 & -0.192 & -0.020 & 0.253 & 0.622 \\
Vino\_122 & 2 & -0.100 & -0.019 & 0.401 & 0.800 \\
\bottomrule
\end{tabular}

\vspace{6pt}
\begin{tabular}{lrrrrrr}
\toprule
Variable & Vino\_159 & Vino\_074 & Vino\_122 & Media C1 & Media C2 & Media C3 \\
\midrule
\texttt{alcohol} & 14.34 & 12.99 & 11.56 & 13.74 & 12.28 & 13.15 \\
\texttt{ac\_malico} & 1.68 & 1.67 & 2.05 & 2.01 & 1.93 & 3.33 \\
\texttt{ceniza} & 2.70 & 2.60 & 3.23 & 2.46 & 2.24 & 2.44 \\
\texttt{alcal\_ceniza} & 25.00 & 30.00 & 28.50 & 17.04 & 20.24 & 21.42 \\
\texttt{magnesio} & 98.00 & 139.00 & 119.00 & 106.34 & 94.55 & 99.31 \\
\texttt{fenoles\_tot} & 2.80 & 3.30 & 3.18 & 2.84 & 2.26 & 1.68 \\
\texttt{flavonoides} & 1.31 & 2.89 & 5.08 & 2.98 & 2.08 & 0.78 \\
\texttt{fenoles\_noflav} & 0.53 & 0.21 & 0.47 & 0.29 & 0.36 & 0.45 \\
\texttt{proantocianinas} & 2.70 & 1.96 & 1.87 & 1.90 & 1.63 & 1.15 \\
\texttt{int\_color} & 13.00 & 3.35 & 6.00 & 5.53 & 3.09 & 7.40 \\
\texttt{tono} & 0.57 & 1.31 & 0.93 & 1.06 & 1.06 & 0.68 \\
\texttt{od280\_od315} & 1.96 & 3.50 & 3.69 & 3.16 & 2.79 & 1.68 \\
\texttt{prolina} & 660.00 & 985.00 & 465.00 & 1115.71 & 519.51 & 629.90 \\
\bottomrule
\end{tabular}
\caption*{Perfil químico de los tres vinos frente al promedio de cada cultivar.}
\end{table}

\begin{itemize}[leftmargin=*]
  \item \textbf{Vino\_122} (cultivar 2): es el que más contribuye a CP3 ($\approx 11.1\,\%$ frente al promedio de 0.56\,\%) y está muy bien representado en ese eje ($\cos^2 \approx 0.75$). Tiene el máximo de la base en \var{ceniza} (3.23) y en \var{flavonoides} (5.08), y una alcalinidad de cenizas muy alta (28.5 frente a 20.2 de su cultivar). Por eso queda arriba del todo en el plano (1,3). También es el mayor atípico en ambas estadísticas $d^2$.
  \item \textbf{Vino\_159} (cultivar 3): segundo mayor contribuyente a CP2 ($\approx 2.8\,\%$), con coordenada muy negativa en ese eje. Tiene la máxima intensidad de color de la base (13.0) y más alcohol (14.3) que el promedio de su cultivar (13.2): es una versión ``extrema'' del perfil alcohol--color.
  \item \textbf{Vino\_074} (cultivar 2): tiene la alcalinidad de cenizas máxima (30.0) y magnesio alto (139 frente a 94.6 de su cultivar). Contribuye $\approx 4.4\,\%$ a CP3 y está entre los atípicos por $d^2_{1i}$. Aunque es del cultivar 2, en el plano (1,2) queda lejos del centro de su grupo.
\end{itemize}

\subsection*{Indicador $I_1$ en $[0,1]$}

$I_1$ reescala la CP1 al intervalo $[0,1]$ usando cuartiles y $3\cdot\mathrm{RIC}$, de modo que un vino extremo no fija por sí solo los límites; por eso los valores observados quedan entre $\approx 0.36$ y $0.65$. El signo de la CP1 en \texttt{sklearn} es \textbf{opuesto} al de la DVS a mano, así que en $I_1$ los valores altos corresponden al lado fenólico. Los 5 máximos son todos del cultivar 1 (Vino\_015, 004, 019, 053, 011) y los 5 mínimos del cultivar 3 (Vino\_147, 138, 137, 172, 156).

\begin{figure}[H]
\centering
\begin{tikzpicture}
\begin{axis}[ejes, width=11cm, height=6.5cm, xlabel={$I_1$ (indicador basado en CP1)}, ylabel={Densidad},
  ymin=0, xmin=0.30, xmax=0.72, legend style={font=\footnotesize, at={(0.02,0.98)}, anchor=north west},
  title={Histograma y densidad del indicador $I_1$}]
\addplot[ybar interval, fill=gray!25, draw=black] coordinates {(0.3564,1.3753) (0.3809,3.2091) (0.4054,5.9597) (0.4300,2.7506) (0.4545,2.5214) (0.4790,4.1259) (0.5035,2.2922) (0.5280,5.2720) (0.5525,4.5844) (0.5770,5.9597) (0.6015,2.2922) (0.6260,0.4584) (0.6505,0)};
\addplot[blue, very thick, smooth] coordinates {(0.3000,0.0200) (0.3035,0.0279) (0.3071,0.0385) (0.3106,0.0524) (0.3141,0.0704) (0.3176,0.0934) (0.3212,0.1223) (0.3247,0.1582) (0.3282,0.2024) (0.3318,0.2558) (0.3353,0.3198) (0.3388,0.3954) (0.3424,0.4838) (0.3459,0.5860) (0.3494,0.7026) (0.3529,0.8343) (0.3565,0.9813) (0.3600,1.1434) (0.3635,1.3200) (0.3671,1.5100) (0.3706,1.7119) (0.3741,1.9233) (0.3776,2.1414) (0.3812,2.3630) (0.3847,2.5842) (0.3882,2.8008) (0.3918,3.0085) (0.3953,3.2030) (0.3988,3.3800) (0.4024,3.5361) (0.4059,3.6682) (0.4094,3.7743) (0.4129,3.8534) (0.4165,3.9055) (0.4200,3.9319) (0.4235,3.9348) (0.4271,3.9174) (0.4306,3.8835) (0.4341,3.8373) (0.4376,3.7832) (0.4412,3.7255) (0.4447,3.6679) (0.4482,3.6137) (0.4518,3.5652) (0.4553,3.5241) (0.4588,3.4911) (0.4624,3.4662) (0.4659,3.4485) (0.4694,3.4371) (0.4729,3.4305) (0.4765,3.4273) (0.4800,3.4262) (0.4835,3.4266) (0.4871,3.4280) (0.4906,3.4307) (0.4941,3.4354) (0.4976,3.4432) (0.5012,3.4557) (0.5047,3.4743) (0.5082,3.5006) (0.5118,3.5358) (0.5153,3.5809) (0.5188,3.6362) (0.5224,3.7016) (0.5259,3.7766) (0.5294,3.8599) (0.5329,3.9500) (0.5365,4.0450) (0.5400,4.1426) (0.5435,4.2402) (0.5471,4.3353) (0.5506,4.4247) (0.5541,4.5056) (0.5576,4.5746) (0.5612,4.6285) (0.5647,4.6643) (0.5682,4.6789) (0.5718,4.6697) (0.5753,4.6345) (0.5788,4.5719) (0.5824,4.4812) (0.5859,4.3624) (0.5894,4.2166) (0.5929,4.0458) (0.5965,3.8527) (0.6000,3.6406) (0.6035,3.4135) (0.6071,3.1754) (0.6106,2.9306) (0.6141,2.6834) (0.6176,2.4377) (0.6212,2.1970) (0.6247,1.9645) (0.6282,1.7429) (0.6318,1.5341) (0.6353,1.3397) (0.6388,1.1606) (0.6424,0.9976) (0.6459,0.8506) (0.6494,0.7194) (0.6529,0.6036) (0.6565,0.5022) (0.6600,0.4144) (0.6635,0.3391) (0.6671,0.2750) (0.6706,0.2212) (0.6741,0.1762) (0.6776,0.1391) (0.6812,0.1088) (0.6847,0.0843) (0.6882,0.0646) (0.6918,0.0490) (0.6953,0.0368) (0.6988,0.0273) (0.7024,0.0200) (0.7059,0.0145) (0.7094,0.0104) (0.7129,0.0073) (0.7165,0.0051) (0.7200,0.0035)};
\addplot[c1, thick, dashed, smooth] coordinates {(0.3000,0.0000) (0.3035,0.0000) (0.3071,0.0000) (0.3106,0.0000) (0.3141,0.0000) (0.3176,0.0000) (0.3212,0.0000) (0.3247,0.0000) (0.3282,0.0000) (0.3318,0.0000) (0.3353,0.0000) (0.3388,0.0000) (0.3424,0.0000) (0.3459,0.0000) (0.3494,0.0000) (0.3529,0.0000) (0.3565,0.0000) (0.3600,0.0000) (0.3635,0.0000) (0.3671,0.0000) (0.3706,0.0000) (0.3741,0.0000) (0.3776,0.0000) (0.3812,0.0000) (0.3847,0.0000) (0.3882,0.0000) (0.3918,0.0000) (0.3953,0.0000) (0.3988,0.0000) (0.4024,0.0000) (0.4059,0.0000) (0.4094,0.0000) (0.4129,0.0000) (0.4165,0.0000) (0.4200,0.0000) (0.4235,0.0000) (0.4271,0.0000) (0.4306,0.0000) (0.4341,0.0000) (0.4376,0.0000) (0.4412,0.0000) (0.4447,0.0000) (0.4482,0.0000) (0.4518,0.0000) (0.4553,0.0000) (0.4588,0.0000) (0.4624,0.0000) (0.4659,0.0001) (0.4694,0.0002) (0.4729,0.0005) (0.4765,0.0015) (0.4800,0.0041) (0.4835,0.0101) (0.4871,0.0231) (0.4906,0.0486) (0.4941,0.0949) (0.4976,0.1726) (0.5012,0.2928) (0.5047,0.4661) (0.5082,0.7001) (0.5118,0.9995) (0.5153,1.3673) (0.5188,1.8056) (0.5224,2.3156) (0.5259,2.8936) (0.5294,3.5260) (0.5329,4.1878) (0.5365,4.8480) (0.5400,5.4821) (0.5435,6.0865) (0.5471,6.6850) (0.5506,7.3232) (0.5541,8.0503) (0.5576,8.8967) (0.5612,9.8559) (0.5647,10.8781) (0.5682,11.8788) (0.5718,12.7579) (0.5753,13.4230) (0.5788,13.8088) (0.5824,13.8857) (0.5859,13.6588) (0.5894,13.1611) (0.5929,12.4466) (0.5965,11.5835) (0.6000,10.6463) (0.6035,9.7035) (0.6071,8.8041) (0.6106,7.9668) (0.6141,7.1806) (0.6176,6.4160) (0.6212,5.6439) (0.6247,4.8527) (0.6282,4.0565) (0.6318,3.2913) (0.6353,2.6013) (0.6388,2.0217) (0.6424,1.5666) (0.6459,1.2261) (0.6494,0.9736) (0.6529,0.7779) (0.6565,0.6141) (0.6600,0.4694) (0.6635,0.3413) (0.6671,0.2331) (0.6706,0.1484) (0.6741,0.0875) (0.6776,0.0477) (0.6812,0.0240) (0.6847,0.0111) (0.6882,0.0047) (0.6918,0.0019) (0.6953,0.0007) (0.6988,0.0002) (0.7024,0.0001) (0.7059,0.0000) (0.7094,0.0000) (0.7129,0.0000) (0.7165,0.0000) (0.7200,0.0000)};
\addplot[c2, thick, dashed, smooth] coordinates {(0.3000,0.0000) (0.3035,0.0000) (0.3071,0.0000) (0.3106,0.0000) (0.3141,0.0000) (0.3176,0.0000) (0.3212,0.0000) (0.3247,0.0000) (0.3282,0.0000) (0.3318,0.0000) (0.3353,0.0000) (0.3388,0.0001) (0.3424,0.0001) (0.3459,0.0003) (0.3494,0.0006) (0.3529,0.0012) (0.3565,0.0024) (0.3600,0.0044) (0.3635,0.0080) (0.3671,0.0140) (0.3706,0.0237) (0.3741,0.0386) (0.3776,0.0608) (0.3812,0.0929) (0.3847,0.1375) (0.3882,0.1977) (0.3918,0.2765) (0.3953,0.3769) (0.3988,0.5018) (0.4024,0.6541) (0.4059,0.8366) (0.4094,1.0520) (0.4129,1.3030) (0.4165,1.5912) (0.4200,1.9173) (0.4235,2.2795) (0.4271,2.6734) (0.4306,3.0917) (0.4341,3.5246) (0.4376,3.9615) (0.4412,4.3922) (0.4447,4.8096) (0.4482,5.2103) (0.4518,5.5957) (0.4553,5.9701) (0.4588,6.3392) (0.4624,6.7065) (0.4659,7.0707) (0.4694,7.4236) (0.4729,7.7505) (0.4765,8.0318) (0.4800,8.2468) (0.4835,8.3776) (0.4871,8.4138) (0.4906,8.3548) (0.4941,8.2109) (0.4976,8.0014) (0.5012,7.7515) (0.5047,7.4879) (0.5082,7.2343) (0.5118,7.0077) (0.5153,6.8163) (0.5188,6.6589) (0.5224,6.5265) (0.5259,6.4051) (0.5294,6.2785) (0.5329,6.1326) (0.5365,5.9570) (0.5400,5.7473) (0.5435,5.5047) (0.5471,5.2352) (0.5506,4.9479) (0.5541,4.6526) (0.5576,4.3581) (0.5612,4.0708) (0.5647,3.7936) (0.5682,3.5265) (0.5718,3.2669) (0.5753,3.0109) (0.5788,2.7546) (0.5824,2.4949) (0.5859,2.2310) (0.5894,1.9643) (0.5929,1.6984) (0.5965,1.4386) (0.6000,1.1914) (0.6035,0.9627) (0.6071,0.7579) (0.6106,0.5805) (0.6141,0.4320) (0.6176,0.3121) (0.6212,0.2187) (0.6247,0.1484) (0.6282,0.0976) (0.6318,0.0621) (0.6353,0.0382) (0.6388,0.0227) (0.6424,0.0131) (0.6459,0.0073) (0.6494,0.0039) (0.6529,0.0020) (0.6565,0.0010) (0.6600,0.0005) (0.6635,0.0002) (0.6671,0.0001) (0.6706,0.0000) (0.6741,0.0000) (0.6776,0.0000) (0.6812,0.0000) (0.6847,0.0000) (0.6882,0.0000) (0.6918,0.0000) (0.6953,0.0000) (0.6988,0.0000) (0.7024,0.0000) (0.7059,0.0000) (0.7094,0.0000) (0.7129,0.0000) (0.7165,0.0000) (0.7200,0.0000)};
\addplot[c3, thick, dashed, smooth] coordinates {(0.3000,0.0000) (0.3035,0.0000) (0.3071,0.0000) (0.3106,0.0000) (0.3141,0.0002) (0.3176,0.0006) (0.3212,0.0021) (0.3247,0.0067) (0.3282,0.0187) (0.3318,0.0466) (0.3353,0.1040) (0.3388,0.2090) (0.3424,0.3803) (0.3459,0.6320) (0.3494,0.9700) (0.3529,1.3937) (0.3565,1.9009) (0.3600,2.4933) (0.3635,3.1744) (0.3671,3.9432) (0.3706,4.7887) (0.3741,5.6933) (0.3776,6.6473) (0.3812,7.6675) (0.3847,8.8012) (0.3882,10.1028) (0.3918,11.5884) (0.3953,13.1969) (0.3988,14.7917) (0.4024,16.2085) (0.4059,17.3154) (0.4094,18.0345) (0.4129,18.3092) (0.4165,18.0563) (0.4200,17.1672) (0.4235,15.5830) (0.4271,13.3917) (0.4306,10.8602) (0.4341,8.3571) (0.4376,6.2126) (0.4412,4.6046) (0.4447,3.5355) (0.4482,2.8890) (0.4518,2.5118) (0.4553,2.2709) (0.4588,2.0727) (0.4624,1.8601) (0.4659,1.6061) (0.4694,1.3104) (0.4729,0.9953) (0.4765,0.6955) (0.4800,0.4430) (0.4835,0.2554) (0.4871,0.1325) (0.4906,0.0617) (0.4941,0.0257) (0.4976,0.0095) (0.5012,0.0032) (0.5047,0.0009) (0.5082,0.0002) (0.5118,0.0001) (0.5153,0.0000) (0.5188,0.0000) (0.5224,0.0000) (0.5259,0.0000) (0.5294,0.0000) (0.5329,0.0000) (0.5365,0.0000) (0.5400,0.0000) (0.5435,0.0000) (0.5471,0.0000) (0.5506,0.0000) (0.5541,0.0000) (0.5576,0.0000) (0.5612,0.0000) (0.5647,0.0000) (0.5682,0.0000) (0.5718,0.0000) (0.5753,0.0000) (0.5788,0.0000) (0.5824,0.0000) (0.5859,0.0000) (0.5894,0.0000) (0.5929,0.0000) (0.5965,0.0000) (0.6000,0.0000) (0.6035,0.0000) (0.6071,0.0000) (0.6106,0.0000) (0.6141,0.0000) (0.6176,0.0000) (0.6212,0.0000) (0.6247,0.0000) (0.6282,0.0000) (0.6318,0.0000) (0.6353,0.0000) (0.6388,0.0000) (0.6424,0.0000) (0.6459,0.0000) (0.6494,0.0000) (0.6529,0.0000) (0.6565,0.0000) (0.6600,0.0000) (0.6635,0.0000) (0.6671,0.0000) (0.6706,0.0000) (0.6741,0.0000) (0.6776,0.0000) (0.6812,0.0000) (0.6847,0.0000) (0.6882,0.0000) (0.6918,0.0000) (0.6953,0.0000) (0.6988,0.0000) (0.7024,0.0000) (0.7059,0.0000) (0.7094,0.0000) (0.7129,0.0000) (0.7165,0.0000) (0.7200,0.0000)};
\legend{Histograma, Densidad global, Cultivar 1, Cultivar 2, Cultivar 3}
\end{axis}
\end{tikzpicture}
\caption{Distribución de $I_1$: la densidad global es ancha porque mezcla tres modos, uno por cultivar.}
\label{fig:I1}
\end{figure}

La densidad global es ancha y casi plana porque mezcla tres grupos: cada cultivar forma su propio modo (cultivar 3 $\approx 0.41$, cultivar 2 $\approx 0.49$, cultivar 1 $\approx 0.58$), y el cultivar 2 se traslapa con los otros dos.

\subsection*{Atípicos multivariados: $d^2_{1i}$ y $d^2_{2i}$}

\[
  d^2_{1i} = \sum_{k=1}^{p_1} \frac{f_{ik}^2}{\lambda_k}, \qquad
  d^2_{2i} = \sum_{k=p-q+1}^{p} \frac{f_{ik}^2}{\lambda_k},
\]
con $p_1 = 3$ (componentes retenidas) y $q = 2$ (últimas componentes, $\lambda_{12}, \lambda_{13} < 0.2$). El umbral es el cuantil empírico del 95\,\%.

\begin{figure}[H]
\centering
\begin{tikzpicture}
\begin{groupplot}[group style={group size=1 by 2, vertical sep=2.2cm},
  width=\textwidth, height=5.2cm, ejes, xmin=0, xmax=179, ymin=0,
  xlabel={Observación}, ylabel={Estadística}]
\nextgroupplot[title={$d^2_{1i}$ -- primeras 3 componentes}]
\addplot[ycomb, black, thin, mark=none] table[x=obs, y=d1]{\tablaD};
\addplot[rojo, dashed, thick, domain=0:179, samples=2] {6.037};
\node[font=\tiny, anchor=south, rotate=0] at (axis cs:15,6.778) {15};
\node[font=\tiny, anchor=south, rotate=0] at (axis cs:26,10.599) {26};
\node[font=\tiny, anchor=south, rotate=0] at (axis cs:60,18.400) {60};
\node[font=\tiny, anchor=south, rotate=0] at (axis cs:74,9.244) {74};
\node[font=\tiny, anchor=south, rotate=0] at (axis cs:116,7.262) {116};
\node[font=\tiny, anchor=south, rotate=0] at (axis cs:122,20.051) {122};
\node[font=\tiny, anchor=south, rotate=0] at (axis cs:123,6.699) {123};
\node[font=\tiny, anchor=south, rotate=0] at (axis cs:128,8.889) {128};
\node[font=\tiny, anchor=south, rotate=0] at (axis cs:159,6.078) {159};
\nextgroupplot[title={$d^2_{2i}$ -- últimas 2 componentes}]
\addplot[ycomb, black, thin, mark=none] table[x=obs, y=d2]{\tablaD};
\addplot[rojo, dashed, thick, domain=0:179, samples=2] {6.185};
\node[font=\tiny, anchor=south] at (axis cs:68,7.126) {68};
\node[font=\tiny, anchor=south] at (axis cs:69,6.688) {69};
\node[font=\tiny, anchor=south] at (axis cs:75,11.244) {75};
\node[font=\tiny, anchor=south] at (axis cs:122,26.863) {122};
\node[font=\tiny, anchor=south] at (axis cs:131,7.318) {131};
\node[font=\tiny, anchor=south] at (axis cs:153,6.391) {153};
\node[font=\tiny, anchor=south] at (axis cs:159,6.438) {159};
\node[font=\tiny, anchor=south] at (axis cs:160,6.394) {160};
\node[font=\tiny, anchor=south] at (axis cs:168,6.991) {168};
\end{groupplot}
\end{tikzpicture}
\caption{Estadísticas de atípicos por observación; la línea roja es el cuantil 95\,\% ($d^2_1$: 6.037; $d^2_2$: 6.185). Se etiqueta el número de cada vino atípico.}
\label{fig:d2}
\end{figure}

\begin{table}[H]
\centering\small
\caption{Vinos atípicos según cada estadística.}
\begin{tabular}{p{3.5cm}p{10cm}}
\toprule
Estadística & Vinos atípicos (cultivar) \\
\midrule
$d^2_{1i}$ ($p_1 = 3$) & 015 (1), 026 (1), 060 (2), 074 (2), 116 (2), \textbf{122} (2), 123 (2), 128 (2), \textbf{159} (3) \\
$d^2_{2i}$ ($q = 2$) & 068 (2), 069 (2), 075 (2), \textbf{122} (2), 131 (3), 153 (3), \textbf{159} (3), 160 (3), 168 (3) \\
\bottomrule
\end{tabular}
\end{table}

Con $\alpha = 0.05$ cada estadística marca 9 vinos ($\approx 5\,\%$ de 178). Solo \textbf{Vino\_122} y \textbf{Vino\_159} aparecen en ambas. Los atípicos de $d^2_{1i}$ son casi todos del cultivar 2 y coinciden con los mayores contribuyentes a CP2/CP3; los de $d^2_{2i}$ son en su mayoría distintos.
\begin{itemize}[leftmargin=*]
  \item $d^2_{1i}$ mide qué tan lejos está un vino del centro \textbf{dentro del subespacio retenido}. Detecta versiones ``exageradas'' del patrón general (mucho color/alcohol o mucha ceniza), que pueden inflar la varianza de las primeras componentes.
  \item $d^2_{2i}$ mide la distancia en las direcciones donde casi no hay variabilidad. Detecta vinos que \textbf{rompen la estructura de correlación}: combinaciones inusuales (p.\,ej., muchos fenoles totales pero pocos flavonoides) aunque cada valor por separado sea normal. Son invisibles en el primer plano.
\end{itemize}

\paragraph{Reflexión.} Los valores atípicos pueden distorsionar los resultados del ACP, porque la varianza y la covarianza son sensibles a valores extremos: unos pocos vinos pueden dominar la dirección de los ejes y llevar a una interpretación incorrecta.

% =====================================================================
\section{Conclusiones y limitaciones}

\paragraph{a) Respuesta a la pregunta de investigación.} Las 13 mediciones químicas se redujeron a tres componentes principales que retienen el 66.5\,\% de la inercia total (CP1: 36.2\,\%, CP2: 19.2\,\%, CP3: 11.1\,\%). Estas dimensiones tienen un sentido químico claro y coherente con la conjetura inicial: CP1 representa la riqueza fenólica y calidad antioxidante; CP2, el perfil cromático y de madurez (color, alcohol, prolina); CP3, la composición mineral. Al proyectar las observaciones en el primer plano factorial, los tres cultivares se separan de manera clara. Las diferencias químicas entre variedades no están dispersas al azar, sino concentradas en estas dimensiones latentes (lideradas por la riqueza fenólica y el color), lo que valida el ACP como herramienta de reducción de dimensionalidad para este conjunto de datos.

\paragraph{b) Información perdida.} Con $k = 3$ se retiene $\approx 66.5\,\%$ de la inercia y se pierde $\approx \mathbf{33.5\,\%}$. La pérdida es aceptable para el objetivo: las 10 componentes descartadas tienen $\lambda < 1$ y la separación de los cultivares ya se logra en el primer plano. Si el objetivo fuera reconstruir con precisión cada medición, convendría retener 5 componentes ($\approx 80\,\%$).

\paragraph{c) Limitaciones.}
\begin{enumerate}[leftmargin=*]
  \item \textbf{Linealidad.} El ACP solo captura relaciones lineales; las variables asimétricas generan nubes en abanico que la correlación de Pearson resume solo parcialmente.
  \item \textbf{Sensibilidad a atípicos y a la escala.} Unos pocos vinos extremos (p.\,ej., Vino\_122, que aporta 11\,\% de la inercia de CP3) influyen en las correlaciones y en la dirección de los ejes. Además, el resultado depende de haber estandarizado.
  \item \textbf{Signo arbitrario de los ejes.} El sentido de cada componente lo decide el algoritmo (DVS, \texttt{eigh}, \texttt{sklearn}); hay que verificarlo antes de afirmar ``más alto = más fenólico''.
  \item \textbf{Ejes con poca inercia.} La CP3 depende casi solo de dos variables, por lo que su interpretación es menos robusta; la CP4 quedó por fuera aunque está cerca del umbral.
  \item \textbf{Representatividad.} Son 178 vinos de una sola región italiana y tres cultivares, analizados hace más de 30 años; las conclusiones no se generalizan a otras regiones, añadas o variedades.
\end{enumerate}

\paragraph{d) Trabajo futuro.}
\begin{itemize}[leftmargin=*]
  \item Comparar con un \textbf{ACP robusto} (MCD o \emph{projection pursuit}) para medir cuánto cambian los ejes sin los atípicos.
  \item Usar \textbf{análisis discriminante lineal} con el cultivar como respuesta, para evaluar si las 3 componentes bastan para predecir la variedad.
  \item Ampliar con datos de \textbf{Wine Quality} (UCI) o variables sensoriales, para relacionar los ejes químicos con la calidad percibida.
\end{itemize}

% =====================================================================
\section*{Referencias}
\addcontentsline{toc}{section}{Referencias}
\begin{itemize}[leftmargin=*, label={}]
  \item Aeberhard, S. y Forina, M. (1992). \emph{Wine} [Dataset]. UCI Machine Learning Repository. \url{https://doi.org/10.24432/C5PC7J}
  \item Pacheco, M. (s.f.). \emph{Análisis Multivariado} [Repositorio de código y material de clase]. \url{https://datastatistic.github.io/multivariate/}
\end{itemize}

\end{document}
